In this section, we study the category of $C\ta$-modules.
Our main theorem states that this category admits an explicit, algebraic description.
Below, we use $\mathrm{IndCoh}(\mathcal{M}_{\mathrm{fg}})$ to denote $\mathrm{Ind}$ of the thick subcategory of $\mathrm{MU}_*\mathrm{MU}$-comodules generated by even shifts of the unit. For more details on this category, see \Cref{dfn:IndCoh}.


\begin{thm} \label{thm:Cta}
  There is an equivalence of presentably symmetric monoidal categories under $\Sp_{C_2, i2}$
  \begin{align} \Mod ( \SHR^{\at} _{i2} ; C\ta ) \simeq \Mod ( \Sp_{C_2, i2}; \uZt) \otimes_{\Z} \mathrm{IndCoh}(\mathcal{M}_{\mathrm{fg}}), \label{eqn:cta} \end{align}
  where $\Sp_{C_2, i2}$ acts on the left-hand-side through $c_{\C/\R}$ and on the left factor of the right-hand-side. %\footnote{In other language, $\IndCoh(\Mfg)$ is the evenly-graded version of Hovey's category $\mathrm{Stable}_{\MU_* \MU}$ of stable $\MU_* \MU$ comodules.} 
	On Picard elements this equivalence takes
  \[ C\ta \otimes \Ss_2^{p,q,w} \mapsto \Sigma^{(p-w)+(q-w)\sigma} \uZt \otimes \omega_{\mathbb{G}/\Mfg} ^{\otimes w}. \]
  %% Under this equivalence the tri-graded sphere $\Ss^{p,q,w}$ is sent to $\Sigma^{(p-w)+(q-w)\sigma} \uZt \otimes \omega_{\mathbb{G}/\Mfg} ^{\otimes w}$.
  As a consequence of this equivalence of categories, there is an isomorphism of tri-graded commutative rings:\footnote{The tensor product may be moved outside the Ext, but then it must be taken in the derived sense.}
  \[ \pi_{p,q,w}^\R(C\ta) \cong \bigoplus_{w+a-s = p} \Ext_{\MU_*\MU}^{s,2w}(\MU_*, \MU_* \otimes \pi_{a + (q-w) \sigma}^{C_2} \uZt). \]
\end{thm}

Over $\C$ the corresponding result is the main theorem of \cite{ctauactalol}.
Upon tensoring with $\Spec(\C)$, our theorem recovers theirs. %Just as in the $\C$ case the significance of this result is that special fiber of the deformation with parameter $\ta$ is algebraic. 

\begin{ntn}
  In the above theorem and throughout the section, we follow the convention of writing $\CC \otimes_R \D$ for the tensor product of two presentable $R$-linear categories (instead of $\CC \otimes_{\Mod_R} \D$).
\end{ntn}

\begin{rmk}
  In the statement of the above theorem, we use the $\Z$-linear structure on $\Mod(\Sp_{C_2, i2};\uZt)$ coming from the symmetric monoidal functor $(\_) \otimes_{\uZ} \uZt : \Mod_\Z \to \Mod(\Sp_{C_2, i2};\uZt)$, where $\_ : \Mod_{\Z} \to \Mod(\Sp_{C_2}; \uZ)$ is the symmetric monoidal functor discussed in \Cref{rmk:underline} below.
\end{rmk}

\begin{rmk}
  We will prove in \Cref{lem:Cta-t-structure} that the category $\Mod ( \Sp_{C_2, i2}; \uZt) \otimes_{\Z} \mathrm{IndCoh}(\mathcal{M}_{\mathrm{fg}})$ admits a $t$-structure with heart the category of comodules in $C_2$-Mackey functors over the Hopf algebroid $(\piu_{*\rho} ^{C_2} (\MU_{\R,2}), \piu_{*\rho}^{C_2} (\MU_{\R, 2} \otimes \MU_{\R, 2}))$.
  
  %As a consequence, we may view $\Mod ( \Sp_{C_2, i2}; \uZt) \otimes_{\Z} \mathrm{IndCoh}(\mathcal{M}_{\mathrm{fg}})$ as a derived 
  We may therefore view any such comodule as an element of $\Mod ( \Sp_{C_2, i2}; \uZt) \otimes_{\Z} \mathrm{IndCoh}(\mathcal{M}_{\mathrm{fg}})$.
\end{rmk}

In fact, we are able to explicitly describe the trigraded homotopy groups of $Y \otimes C\ta$ for a more general collection of $Y$ than just trigraded spheres.

\begin{dfn}
  We say that $X \in \Sp_{C_2, i2}$ is \emph{$\MU_{\R,2}$-projective}
  if $\MU_{\R,2} \otimes X$ is a retract (as an $\MU_{\R,2}$-module) of
  $\bigoplus_\alpha \Sigma^{n_\alpha \rho} \MU_{\R,2}$
  for some set of integers $n_\alpha$.
\end{dfn}

Examples of $\MU_{\R,2}$-projective $X$ include $\Ss_2 ^0$ and $\MU_{\R,2}$, as well as any object of $\Sp_{C_2, i2}$ built out of cells of the form $\Ss_2^{n\rho}$.

\begin{thm} \label{thm:Cta-homotopy}
  If $X \in \Sp_{C_2, i2}$ is $\MUR$-projective,
  then under the equivalence of \Cref{thm:Cta}  
  $C\ta \otimes \Gamma_*(X)$ corresponds to the $\piu_{*\rho} ^{C_2} (\MU_{\R,2} \otimes \MU_{\R,2})$-comodule
  $\piu_{*\rho} ^{C_2} (\MU_{\R,2} \otimes X)$.
  %\footnote{We will prove in \Cref{} that the category $\Mod ( \Sp_{C_2, i2}; \uZt) \otimes_{\Z} \mathrm{IndCoh}(\mathcal{M}_{\mathrm{fg}})$ admits a $t$-structure with heart the category of comodules in $C_2$-Mackey functors with over the Hopf algebroid $(\piu_{*\rho} ^{C_2} (\MU_{\R,2}), \piu_{*\rho}^{C_2} (\MU_{\R, 2} \otimes \MU_{\R, 2}))$ }

  Moreover, there is an isomorphism of trigraded groups,
  \[ \pi_{p,q,w} ^\R (\Gamma_* (X) \otimes C\ta) \cong \bigoplus_{w+a-s = p} \Ext_{(\MU_2)_* \MU_2}^{s,2w}((\MU_2)_*, (\MU_2)_{*} (\Phi^{e} (X)) \otimes_{\Z_2} \pi_{a + (q-w) \sigma}^{C_2} \uZt), \]
  compatible with the $\pi^{\R} _{p,q,w}C\ta$-module structure in the expected way.
  
 % satisfies $(\underline{\MU}_{\R,2})_{n\rho -1} X = 0$ for all $n \in \Z$, then $\Gamma_* (X)$ corresponds to the $(\underline{\MU}_{\R,2})_{*\rho} \MU_{\R,2}$-module $(\underline{\MU}_{\R,2}) _{*\rho} X$ under the above equivalence.
%
%  Suppose further that the restriction map is an isomorphism on $(\underline{\MU}_{\R,2}) _{*\rho} X$.
%  Then there is an isomorphism of trigraded groups:
%  \[ \pi_{p,q,w} ^\R (\Gamma_* (X) \otimes C\ta) \cong \bigoplus_{w+a-s = p} \Ext_{\MU_*\MU}^{s,2w}(\MU_*, \MU_{*} (U(X))) \otimes \pi_{a + (q-w) \sigma}^{C_2} \uZt, \]
%  where the tensor product is derived.
\end{thm}

%% \begin{rmk}
%%   We will call $C_2$-spectra $X$ satisfying the hypotheses of \Cref{thm:Cta-homotopy} $\MUR$-projective.
%%   Examples of $\MUR$-projective $C_2$-spectra include $\Ss^0, \MUR, \ku_{\R}, \ko_{C_2}$, as well as any $C_2$-spectrum built out of cells of the form $\Ss^{n\rho}$.
%%   If it were known to exist, $\tmf_{C_2}$ would also be $\MUR$-projective.
%% \end{rmk}

\todo{There's a commented rmk here that should be exported to somewhere else.}
%% \begin{rmk}
%%   For a $\MUR$-projective $C_2$-spectrum $X$, the $\ta$-Bockstein spectral sequence for $\Gamma_* (X)$ is an interesting mix of the slice spectral sequence and the classical Adams-Novikov spectral sequence, as evidenced by the computation of the $\mathrm{E}_1$-term given in \Cref{thm:Cta-homotopy}.
%% \end{rmk}

The proofs of Theorems \ref{thm:Cta} and \ref{thm:Cta-homotopy} will be quite long, and correctly handling the symmetric monoidal structures involved requires us to take a rather circuitous route. For this reason, before proceeding we provide a sketch of the argument.

The proof of the equivalence
\[ \Mod ( \SHR^{\at} _{i2} ; C\ta ) \simeq \Mod ( \Sp_{C_2, i2}; \uZt) \otimes_{\Z} \mathrm{IndCoh}(\mathcal{M}_{\mathrm{fg}}) \]
has three main steps and one significant subtlety.
First, we % pick an auxiliary category $\CC$ and
produce symmetric monoidal functors into each side of \Cref{eqn:cta} from $\Sp_{C_2,i2}^{\Fil}$.
%% of the categories of interest.
Second, we show that there are commutative algebras $R_1$ and $R_2$ in $\Sp_{C_2, i2}^{\Fil}$ such that the left-hand-side is equivalent to $\Mod ( \Sp_{C_2, i2}^{\Fil} ; R_1)$ and the right-hand-side is equivalent to $\Mod ( \Sp_{C_2, i2}^{\Fil} ; R_2)$.
Third, we examine $R_1$ and $R_2$ directly and find that they are in fact equivalent.

%% The auxiliary category $\CC$ we use is $i2$-complete filtered $C_2$-spectra.
The functor into $C\ta$-modules from $\Sp_{C_2,i2}^{\Fil}$ is the composite of the functor $i_*$ from \Cref{sec:top} with $C\ta \otimes -$.
The functor into the right-hand-side is more delicate to construct. A first guess would be to use the equivalence $\Sp_{C_2,i2}^{\Fil} \cong \Sp_{C_2,i2} \otimes \Sp^{\Fil}$ and tensor the following pair of maps:
\[ \Sp_{C_2, i2} \to \Mod ( \Sp_{C_2, i2} ; \uZt ), \]
\[ \Sp^{\Fil} \to \Mod ( \Sp^{\Fil} ; \Z) \xrightarrow{\Gr} \Mod ( \Sp^{\Gr} ; \Z ) \to \mathrm{IndCoh}(\mathcal{M}_{\mathrm{fg}}) . \]
In fact, this does not produce the correct functor; in $C\ta$-modules the periodicity class $v_1$ lives in $C_2$-degree $\rho$, but the given tensor product functor puts $v_1$ in $C_2$-degree 0. In order to fix this, we twist by a functor
\[ \Mod (\Sp_{C_2,i2}^{\Gr}; \uZt) \xrightarrow{\mathrm{tw}^{\rho}} \Mod (\Sp_{C_2,i2}^{\Gr}; \uZt)\]
%\[ \Sp_{C_2,i2}^{\Fil} \to \Mod (\Sp_{C_2,i2}^{\Gr}; \uZt) \xrightarrow{\mathrm{tw}^{\rho}} \Mod (\Sp_{C_2,i2}^{\Gr}; \uZt) \to \Mod (\Sp_{C_2, i2}^{\Gr}; \uZt) \otimes_\Z \mathrm{IndCoh}(\mathcal{M}_{\mathrm{fg}})\]
that has the effect of tensoring with $\Ss^{n\rho}$ on the $n^{\mathrm{th}}$ graded piece.
The presence of this twist, and its interaction with the symmetric monoidal structure, is the main subtlety of the argument.
\todo{The absence of this twist in the GIKR argument indicates a lack of rigor.}

The second step is straightforward, and demonstrates the power of higher algebra in proving ``Koszul duality'' statements as a corollary of Barr--Beck--Lurie monadicity. At this point we have two commutative algebras $R_1$ and $R_2$, and all we need to do is show they are equivalent.

The third step relies on special properties of $R_1$ and $R_2$.
Both commutative algebras come equippped with a preferred presentation as a totalization of a cosimplicial diagram of commutative algebras. What we do is show that these cosimplicial diagrams are levelwise equivalent. Each cosimplicial diagram takes values in in the heart of a $t$-structure, and in particular is determined by $1$-categorical data. Thus, it is genuinely possible to write down a comparison map by hand and check that it is an equivalence, without concern for the potential infinity of higher coherence data.

%Since this proof sketch doesn't track too closely with the division of the material into subsections, we also provide an outline of the section below. 

\begin{enumerate}
  \item In \Cref{subsec:cats}, we review the categories which appear in the right-hand-side of \Cref{thm:Cta}. We also construct a $t$-structure on the right-hand-side of \Cref{thm:Cta} which will play an important role in our argument.
\item In \Cref{subsec:twists}, we collect some material on twisted $t$-structures on categories of graded and filtered objects, as well as on twisting isomorphisms between them.
\item In \Cref{subsec:slice}, we relate twisted $t$-structures to the slice filtration in $C_2$-equivariant homotopy theory. This finishes the construction of the comparison functors.
\item In \Cref{subsec:reds}, we use Koszul duality to reduce the proof of Theorems \ref{thm:Cta} and \ref{thm:Cta-homotopy} to understanding a specific pair of commutative algebras.
\item In \Cref{subsec:main-lem}, we prove \Cref{lem:main-lem}, which is the key result that allows us to compare the commutative algebras $R_1$ and $R_2$ above. Using this lemma, we then complete the proof of Theorems \ref{thm:Cta} and \ref{thm:Cta-homotopy}.
\end{enumerate}
\todo{This list still needs some attention}

%\begin{rmk}
  %The first condition in the statement of \Cref{thm:Cta-homotopy} corresponds to the condition that $\MU_{\R,2} \otimes X$ be \emph{even} in the sense of \cite[Definition 3.1]{HillMeier}.
  %The second condition corresponds to the further condition that $\MU_{\R,2} \otimes X$ be \emph{strongly even}.
%\end{rmk}

%\begin{rmk}
  %For nice $X$, the restriction map $(\MU_{\R,2}) _{*\rho} X \to (\MU_{2})_{2*} X$ is an equivalence. For example, the is the case for $X = \MU_{\R,2}, \ku_{\R,2}, \ko_{C_2,2}$. In this case, the trigraded homotopy of $\Gamma_* (X) \otimes C\ta$ is determined by taking the derived tensor product of the $\mathrm{E}_2$-page of the Adam-Novikov spectral sequence for the underlying spectrum of $X$ with the $RO(C_2)$-graded homotopy of $\uZt$.
%\end{rmk}

\subsection{Categories of interest} \label{subsec:cats}\ 

In this subsection we set notation for working in the various categories of interest.
%Our goal here is to provide a gentle introduction so that reader interested in working with $C\ta$-modules has good computational control over the category.
%This means digesting the category on the right-hand-side of \Cref{thm:Cta}
In particular, we aim for a practical understanding of the category
\[ \Mod ( \Sp_{C_2, i2}; \uZt) \otimes_{\Z} \mathrm{IndCoh}(\mathcal{M}_{\mathrm{fg}}) \]
that appears on the right-hand-side of \Cref{thm:Cta}. Proceeding from inside out,
we start by fixing notation for the category of $\Z$-modules:

\begin{dfn}
  Let $\Abh$ denote the abelian category of discrete abelian groups.
  Let $\Ab$ denote the category of $\Z$-modules with its standard $t$-structure,
  so that $\Abh$ is the heart of $\Ab$.
\end{dfn}

%From here this subsection now breaks into three parts.
%First, we discuss the relevant categories of equivariant objects.
%Next, we discuss the relevant category of $\MU_*\MU$-comodules, or equivalently sheaves on the moduli stack of formal groups $\Mfg$.
%Finally, we explain how to combine these two pieces.
%Thankfully, the equivariant and comodule aspects are only weakly coupled so we will be able to give simple descriptions.

%% We begin by recalling the definitions of the relevant equivariant categories. These include (discrete) Mackey functors for $C_2$ and modules over $\uZt$ in Mackey functors.
%% Let us begin by taking the time to carefully set up the various players in the story. We start by discussing some categories of Mackey functors and functors between them.

Next, we describe the $C_2$-equivariant analog of $\Ab$:

\begin{dfn} \label{dfn:mackey}
  Let $M(C_2)^{\heartsuit}$ denote the abelian category of discrete Mackey functors for the group $C_2$,
  which may be explicitly described as follows.
  A Mackey functor $B \in M(C_2) ^{\heartsuit}$ consists of the following:
  \begin{itemize}
    \item abelian groups $B(C_2)$ and $B(\ast)$,
    \item an involution $\sigma : B(C_2) \to B(C_2)$,
    \item homomorphisms $r : B(\ast) \to B(C_2)$ and $t : B(C_2) \to B(\ast)$,
     which satisfy the relations $\sigma \circ r = r$, $t \circ \sigma = t$ and $r \circ t = 1 + \sigma$.
  \end{itemize}
  We will somtimes write $[C_2]$ for the endomorphism $t \circ r$ of $B(*)$.
  We let $M(C_2)$ denote the unbounded derived category of $M(C_2)^{\heartsuit}$ equipped with its standard $t$-structure.\footnote{As defined, for example, in \cite[Definition 1.3.5.8]{HA}.}
\end{dfn}

In our setting, $M(C_2)$ arises naturally because it is equivalent to the category of $\piu_0 ^{C_2} \Ss$-modules in $\Sp_{C_2}$ \cite[Theorem 5.10]{SpectrumMackey}.
This description equips $M(C_2)$ with a symmetric monoidal structure compatible with the $t$-structure. The induced symmetric monoidal structure on $M(C_2)^{\heartsuit}$ is known as the box product and admits an explicit description as in \cite{CpBox}.
\todo{Probably there's a better cite for this, but this one has a nice modern proof.}
\todo{rwb: Its always unnerved me that I don't know an explicit description of the monoidal structure on Mackey functors.}

\begin{exm}
  Given a discrete abelian group $A$, we may consider the constant Mackey functor $\underline{A} \in M(C_2)^{\heartsuit}$. This Mackey functor is determined by $\underline{A} (C_2) = \underline{A} (\ast) = A$, $r = \sigma = \mathrm{id}_{A}$ and $t=2 \cdot \mathrm{id}_{A}$.
  This construction provides an exact colimit-preserving functor
  \[\_ : \Abh \to M(C_2)^{\heartsuit}.\]
  Examining the explicit description of the box product in \cite{CpBox}, we find that $\_$ acquires the structure of a (non-unital) tensor-product preserving lax symmetric monoidal functor.
  In particular, we obtain a commutative algebra $\uZ$ in $M(C_2)^{\heartsuit}$.
\end{exm}
\todo{rwb: Given we describe nothing about the monoidal structure, why is this functor lax-mon?}

\begin{dfn}
  Let $\underline{\mathrm{Ab}}$ denote the the symmetric monoidal category of $\uZ$-modules in $M(C_2)$ (or, equivalently, $\Sp_{C_2}$), equipped with its standard $t$-structure. This is equivalent to the derived category of $\Mod(M(C_2)^{\heartsuit}; \uZ)$; in particular, we have $\uAbh = \Mod(M(C_2) ^{\heartsuit}; \uZ)$.
\end{dfn}

\begin{rmk} \label{rmk:underline}
  Since $\Z$ is the unit of $\Abh$, we learn that $\_$ factors through an exact colimit-preserving symmetric monoidal functor% \footnote{A priori this functor is only lax symmetric monoidal, but exact colimit-preserving lax symmetric monoidal functors out of the category of abelian groups which preserve the unit are automatically symmetric monoidal.}
  \[\_ : \Abh \to \uAbh.\]
  Since $\_$ is exact, colimit-preserving and symmetric monoidal it extends uniquely to a colimit-preserving symmetric monoidal functor $\_ : \Ab \to \uAb$ of derived categories. Moreover, the pair of functors $A \to A(*)$ and $A \to A(C_2)$ extend to limit-preserving functors on the level of the derived category. Since this pair of functors is jointly conservative, we can use the fact that $\underline{A}(*) \cong A$ and $\underline{A}(C_2) \cong A$ to conclude that $\_$ preserves limits as well.
 
%% \todo{Is it symmetric monoidal? If so, reason/cite? rwb: Apply BBL to the algebra map Sp --> uAb and you obtain the same symm mon functor in from Ab. Is this worth stating? rwb (later): yes, b/c nobodies written down a lax tensor product of infinity categories we need the functor to have been monoidal.}
\end{rmk}

Before proceeding to explicitly describe $\IndCoh(\Mfg)$, we record a few lemmas for later use. 

\begin{lem} \label{lem:const-mackey}
  The functor $\_ : \Abh \to \uAbh$ is fully faithful and left adjoint to the functor
  $B \mapsto B(\ast)$.
  Its essential image is therefore a coreflective subcategory of $\uAbh$, with coreflector given by the counit map $\underline{B(\ast)} \to B$.
  In particular, for any $B$ in the essential image of $\_$, there are canonical isomorphisms
  \[B \cong \underline{B(\ast)} \cong \underline{B(C_2)}.\]
\end{lem}

\begin{proof} Clear.
\end{proof}

\begin{lem} \label{lem:uabh-image}
  The forgetful functor $\uAbh \to M(C_2)^{\heartsuit}$ is fully faithful, with image spanned by the Mackey functors $B$ for which $[C_2]=2$ as endomorphisms of $B(\ast)$.
\end{lem}

\begin{proof}
  The unit of $M(C_2)^{\heartsuit}$ is the Burnside Mackey functor $\underline{\mathcal{A}}$, whose values are given by $\underline{\mathcal{A}} (C_2) = \Z$ and $\underline{\mathcal{A}} (\ast) = \Z[[C_2]]/([C_2]^2 - 2[C_2])$.
  As a commutative algebra in $M(C_2)^{\heartsuit}$, $\uZ$ is the quotient of $\underline{\mathcal{A}}$ by the relation $[C_2]=2$.
  The result follows.
\end{proof}

The following definition and construction sum up what we need to know about categories of quasicoherent and ind-coherent sheaves on $\Mfg$:

\begin{dfn} \label{dfn:IndCoh}
  We let $\Mfg$ denote the moduli stack of formal groups and let $\QCoh (\Mfg)$ denote the category of quasicoherent sheaves on $\Mfg$.
  It is equivalent to the derived category of evenly graded $\MU_* \MU$-comodules \cite[Remarks 2.38 and 3.14]{GoerssMfg}.
%
  %Let $\QCoh(\Mfg)_{i2}$ denote the category of modules over the $2$-completion of the unit.

  Let $\mathcal{D} \subset \QCoh(\Mfg)$ denote the thick subcategory generated by the sheaves $\omega_{\mathbb{G}/\Mfg} ^{\otimes k}$ for $k \in \Z$, where $\omega_{\mathbb{G} / \Mfg}$ is the sheaf of invariant differentials on the universal formal group $\mathbb{G} / \Mfg$. We define $\IndCoh(\Mfg) \coloneqq \Ind(\mathcal{D})$.
  %It follows from \cite[Remark 4.30]{LocDual} that $\IndCoh(\Mfg)$ is equivalent to .
  The category $\IndCoh(\Mfg)$ is equivalent to Hovey's category of evenly graded stable comodules $\mathrm{Stable}_{\MU_* \MU}$, c.f. \cite[Remark 4.30]{LocDual}.
  %
  %Finally, we let $\IndCoh(\Mfg)_{i2}$ denote the category of modules over the $2$-completion of the unit.

  Both $\QCoh(\Mfg)$ and $\IndCoh(\Mfg)$ are naturally stable presentably symmetric monoidal categories and come equipped with compatible $t$-structures whose hearts are equivalent to the $1$-category of $\MU_* \MU$-comodules. \todo{Come back to this and remark that D = Perf and justify that Indcoh = Indperf here.}
\end{dfn}

\begin{cnstr}
  Since $\QCoh(\Mfg)$ is equivalent to the derived category of a Grothendieck abelian category, it is equipped with the structure of a colimit-preserving symmetric monoidal functor $\Ab \to \QCoh(\Mfg)$.
  Letting $\Ab ^{\mathrm{fin}}$ denote the full subcategory of compact objects, this restricts to an exact symmetric monoidal functor $\Ab ^{\mathrm{fin}} \to \D$, where $\D$ is defined as in \Cref{dfn:IndCoh}.
  Taking $\Ind$ and using the fact that $\Ab$ is compactly generated, we obtain a colimit-preserving symmetric monoidal functor
  \[\Ab \simeq \Ind(\Ab ^{\mathrm{fin}}) \to \Ind(\D) = \IndCoh(\Mfg).\]

  Using this functor and the functor $\Ab \xrightarrow{\_} \uAb \xrightarrow{- \otimes_{\uZ} \uZt} \uAb_{i2}$, we are able to define the stable presentably symmetric monoidal category
  \[\uAb_{i2} \otimes_{\Z} \IndCoh(\Mfg),\]
  which appeared in \Cref{thm:Cta} with slightly different notation.
\end{cnstr}

The core of our understanding of $\uAb_{i2} \otimes_{\Z} \IndCoh(\Mfg)$ rests on the following lemma, where we put a $t$-structure on this category whose heart we can describe explicitly.

\begin{lem} \label{lem:Cta-t-structure}
  There exists a $t$-structure on $\uAb_{i2} \otimes_{\Ab} \IndCoh(\Mfg)$ with the following properties:
  \begin{enumerate}
    \item It is compatible with the symmetric monoidal structure, accessible, compatible with filtered colimits and right complete.
    \item The natural functor
      \[\uAb_{i2, \geq 0} \otimes _{\Ab_{\geq 0}} \IndCoh(\Mfg)_{\geq 0} \to \uAb_{i2} \otimes_{\Ab} \IndCoh(\Mfg)\]
      is fully faithful with essential image the connective objects.
    \item The natural functor
      \[\uAb_{i2} ^{\heartsuit} \otimes_{\Ab^{\heartsuit}} \IndCoh(\Mfg)^{\heartsuit} \to \uAb_{i2} \otimes_{\Ab} \IndCoh(\Mfg)\]
      is fully faithful with essential image the heart of $\uAb_{i2} \otimes_{\Ab} \IndCoh(\Mfg)$.
  \end{enumerate}
\end{lem}

Using the explicit description of $\uAbh$ given by \Cref{dfn:mackey} and \Cref{lem:uabh-image} and the fact that $\IndCoh(\Mfg)^{\heartsuit}$ is equivalent to the category of evenly-graded $\MU_* \MU$-comodules, we obtain the following description of
\[\uAb_{i2}^{\heartsuit} \otimes_{\Ab^{\heartsuit}} \IndCoh(\Mfg)^{\heartsuit}:\]
this is the category of pairs of $(\MU_2)_*\MU_2$-comodules $A(C_2)$ and $A(\ast)$ equipped with following structures (all of which are compatible with the comodule structure):
\begin{itemize}
  \item maps $r : A(\ast) \to A(C_2)$ and $t : A(C_2) \to A(\ast)$, along with an involution $\sigma : A(C_2) \to A(C_2)$,
  \item which satisfy the relations $\sigma \circ r = r$, $t \circ \sigma = t$, $t \circ r = 2$ and $r \circ t = 1 + \sigma$.
\end{itemize}

In other words, $\uAb^{\heartsuit}_{i2} \otimes_{\Ab^{\heartsuit}} \IndCoh(\Mfg)^{\heartsuit}$ is the category of evenly graded comodules over the Hopf algebroid $(\underline{(\MU_2)_{*}}, \underline{(\MU_2)_* \MU_2})$ in graded $C_2$-Mackey functors.

\begin{rmk}
  Using the fact that $$((\underline{\MU_{\R,2}})_{*\rho}, (\underline{\MU_{\R,2}})_{*\rho} \MU_{\R,2}) \cong (\underline{(\MU_2)_{2*}}, \underline{(\MU_2)_{2*} \MU_2}),$$ see \cite[Theorems 2.25 and 2.28]{HuKriz}, this category may equally well be described as the category of graded comodules over the Hopf algebroid $((\underline{\MU_{\R,2}})_{*\rho}, \underline({\MU_{\R,2}})_{*\rho} \MU_{\R,2})$ in graded $C_2$-Mackey functors.
\end{rmk}

The proof of \Cref{lem:Cta-t-structure} is not difficult, but it will rely on material from \cite[Appendix C]{SAG} which we presently recount. We begin with a simple method for producing $t$-structures. 

\begin{cnstr}
  Given a pointed presentable category $\CC$, we let 
  $\Sp \otimes \CC \simeq \Sp(\CC)$
  %% = \varprojlim \left( \dots \xrightarrow{\Omega} \CC \xrightarrow{\Omega} \CC \xrightarrow{\Omega} \CC \right)\]
  denote the category of spectrum objects in $\CC$. It is a stable presentable category.

  The category $\Sp \otimes \CC$ admits a $t$-structure with $(\Sp \otimes \CC)_{\geq 0}$ equal to the essential image of the functor $\Sigma^{\infty} : \CC \to \Sp \otimes \CC$.
  This $t$-structure is accessible and right complete.

  Note that if $\CC$ is endowed with the structure of a presentably symmetric monoidal category, then $\Sp \otimes \CC$ naturally inherits this structure and the $t$-structure defined above is compatible with it.
\end{cnstr}

In fact, the following lemma, which follows from \cite[Remark C.3.1.5]{SAG}, demonstrates that this construction is the universal way to produce a (well-behaved) $t$-structure on a presentable stable category. 

\begin{lem} \label{lem:conn-enough}
  Suppose that $(\CC, \CC_{\geq 0})$ is a stable presentably symmetric monoidal category with compatible $t$-structure. Suppose that the $t$-structure is accessible, compatible with filtered colimits and right complete.

  Then there is a natural equivalence of symmetric monoidal categories with compatible $t$-structure
  \[(\CC, \CC_{\geq 0}) \simeq (\Sp \otimes \CC_{\geq 0}, (\Sp \otimes \CC_{\geq 0})_{\geq 0}).\]
\end{lem}

%% \begin{cor} \label{cor:is-sp}
%%   There are natural equivalences of symmetric monoidal categories with compatible $t$-structures:
%%   \[(\Ab,\Ab_{\geq 0}) \simeq (\Sp(\Ab_{\geq 0}), \Sp(\Ab_{\geq 0})_{\geq 0}),\]
%%   \[(\uAb,\uAb_{\geq 0}) \simeq (\Sp(\uAb_{\geq 0}), \Sp(\uAb_{\geq 0})_{\geq 0}),\]
%%   \[(\IndCoh(\Mfg), \IndCoh(\Mfg)_{\geq 0}) \simeq (\Sp(\IndCoh(\Mfg)_{\geq 0}), \Sp(\IndCoh(\Mfg)_{\geq 0})_{\geq 0}).\]
%% \end{cor}

In conclusion, we may as well work with the categories of connective objects. For $\Ab, \uAb_{i2}$, and $\IndCoh(\Mfg)$, these are all Grothendieck prestable categories in the sense of \cite[Definition C.1.4.2]{SAG}.

\begin{rec} \label{rec:groth-ab}
  We let $\Groth_{\infty}$ denote the category of Grothendieck prestable categories \cite[Definition C.3.0.5]{SAG}. By \cite[Theorem C.4.2.1]{SAG}, $\Groth_{\infty}$ is a symmetric monoidal category with tensor product given by the usual tensor product of presentable categories and unit $\Sp_{\geq 0}$.

  Moreover, let $\Groth_{1}$ denote the category of Grothendieck abelian $1$-categories. It is a symmetric monoidal category with tensor product given by the usual tensor product of presentable categories and unit $\Abh$ \cite[Corollary C.5.4.19]{SAG}. There is a functor $\tau_{\leq 0} : \Groth_{\infty} \to \Groth_{1}$ sending a Grothendieck prestable category to its subcategory of discrete objects.
  By \cite[Remark C.5.4.20]{SAG}, the functor $\tau_{\leq 0}$ is symmetric monoidal.

  We summarize this in the following span of symmetric monoidal categories 
  \begin{center}
    \begin{tikzcd}
      & \Groth_{\infty} \ar[dl, "\Sp \otimes -"] \ar[dr, "\tau_{\leq 0}"] \\
      \mathrm{Pr}^{L,\mathrm{Stab}} & & \Groth_{1}.
    \end{tikzcd}
  \end{center}
\end{rec}

\begin{proof}[Proof (of \Cref{lem:Cta-t-structure}).]
  Using \Cref{lem:conn-enough} and \cite[Remark C.4.2.3]{SAG}, we have
  \[\Sp \otimes \left(\uAb_{i2, \geq 0} \otimes_{\Ab_{\geq 0}} \IndCoh(\Mfg)_{\geq 0} \right) \simeq \uAb_{i2} \otimes_{\Ab} \IndCoh(\Mfg).\]
  Moreover, since $\uAb_{i2, \geq 0} \otimes_{\Ab_{\geq 0}} \IndCoh(\Mfg)_{\geq 0}$ is prestable, the $\Sigma^{\infty}$ functor identifies it with $\left(\Sp \otimes \left(\uAb_{i2, \geq 0} \otimes_{\Ab_{\geq 0}} \IndCoh(\Mfg)_{\geq 0} \right)\right)_{\geq 0}$.

  This immediately implies the first two properties, with the exception of compatibility with filtered colimits. This follows from the fact that $\uAb_{i2, \geq 0} \otimes_{\Ab_{\geq 0}} \IndCoh(\Mfg)_{\geq 0}$ is compactly generated, which holds becuase each one of $\uAb_{i2, \geq 0}$, $\Ab_{\geq 0}$ and $\IndCoh(\Mfg)_{\geq 0}$ is compactly generated \cite[Corollary C.6.2.3]{SAG}.

  The third property follows from the fact that $\CC^{\heartsuit} = \tau_{\leq 0} \left( \CC_{\geq 0} \right)$ and the fact that $\tau_{\leq 0}: \Groth_{\infty} \to \Groth_1$ is symmetric monoidal.
\end{proof}

%\begin{dfn}
%  Let $\CC$ and $\D$ denote stable presentably symmetric monoidal $\infty$-categories equipped with $t$-strutures generated by connective objects $\CC_{\geq 0}$ and $\D_{\geq 0}$. We make the following additional assumptions on the $t$-structure:
%  \begin{itemize}
%    \item We have $\o_\CC \in \CC_{\geq 0}$ and $\o_{\D} \in \D_{\geq 0}$.
%    \item The full subcategories $\CC_{\geq 0}$ and $\D_{\geq 0}$ are closed under the tensor product.
%    \item The full subcategories $\CC_{\geq 0}$ and $\D_{\geq 0}$ of coconnective objects are closed under filtered colimits.
%  \end{itemize}
%\end{dfn}
%
%Remark C.5.4.20
%
%Finally, we define a $t$-structure on $\uAb \otimes_\Z \IndCoh(\Mfg)$ by setting $\left( \uAb \otimes_\Z \IndCoh(\Mfg) \right) _{\geq 0}$ to be the category generated under colimits and extensions by the image of the restriction of the functor $m : \uAb \times \IndCoh(\Mfg) \to \uAb \otimes_{\Z} \IndCoh(\Mfg)$ to $\uAb_{\geq 0} \times \IndCoh(\Mfg)_{\geq 0}$.
%
%
%We may now make sense of the tensor product on the right-hand side of the isomorphism in \Cref{thm:Cta}. By the description of $\uAb$ in \Cref{prop:uA-pre}, we see that there is an equivalence
%\[\Mod(\Sp_{C_2}; \uZ) \otimes_{\Z} \IndCoh(\Mfg) = \uAb \otimes_{\Z} \IndCoh(\Mfg) \simeq \Fun^{\times} (\uA, \IndCoh(\Mfg)).\]
%



%
%The usual $t$-structure on $\IndCoh(\Mfg)$ induces a $t$-structure on $\Fun^{\times} (\uA, \IndCoh(\Mfg))$ with heart $\Fun^{\times} (\uA, \IndCoh(\Mfg)^{\heartsuit})$. Unpacking the definitions, we find that an object of this heart consists \todo{up to p-completion nonsense...} of a pair of $\MU_*\MU$-comodules $A(C_2)$ and $A(\ast)$ equipped with following strucutre (all of which is compatible with the comodule structure),
%\begin{itemize}
  %\item an involution, $\sigma$, on $A(C_2)$,
  %\item restriction and transfer maps, $r: A(\ast) \to A(C_2)$ and $t: A(C_2) \to A(\ast)$,
%\end{itemize}
%subject to compatabilities $r \sigma = r$, $\sigma t = t$, $rt = 2$ and $tr = 1 + \sigma$.
%
%
%
%Now we begin our task of understanding the tensor product on the right-hand-side of the isomorphism in the statement of [thm]. Since $\uAb$ can be described as product preserving presheaves of $\Z$-modules on $\uA$ we can correspondingly describe $\Mod ( \Sp_{C_2}; \uZ) \otimes_{\Z} \mathrm{IndCoh}(\mathcal{M}_{\mathrm{fg}})$ as product preserving presheaves on $\uA$ with values in $\mathrm{IndCoh}(\mathcal{M}_{\mathrm{fg}})$.  
%
%By declaring that external tensors of connective objects are connective we obtain a $t$-structure on $\Mod ( \Sp_{C_2}; \uZ) \otimes_{\Z} \mathrm{IndCoh}(\mathcal{M}_{\mathrm{fg}})$. The sheaf description gives us a simple description of the heart of this $t$-structure. 
%An object of $(\uAb \otimes_{\Ab} \mathrm{IndCoh}(\mathcal{M}_{\mathrm{fg}}))^\heartsuit$
%consists of a pair of $\MU_*\MU$-comodules $A_e$ and $A_{C_2}$ equipped with following strucutre (all of which is compatible with the comodule structure),
%\begin{itemize}
%\item an involution, $\sigma$, on $A_e$,
%\item restriction and transfer maps, $r: A_{C_2} \to A_e$ and $t: A_e \to A_{C_2}$,
%\end{itemize}
%subject to compatabilities $r \sigma = r$, $\sigma t = t$, $rt = 2$ and $tr = 1 + \sigma$.
%
%Using this description we can conclude that the statement about the trigraded homotopy ring in [thm] follows from equivalence of categories given earlier in [thm]. \todo{Actually prove this and come back to it later.}
%We will prove the remainder of [thm] over the course of the next two subsections.
%In [subsec], we will use general nonsense to produce a comparison functor between the two categories in question and reduce [thm] to a single, more technical lemma.
%In [subsec], we will use the special properties of $t$-structures to prove this lemma.
%\todo{Update this at the end to take into account the fact that there are two theorems now.}
%
%% \begin{rmk}\todo{Move this elsewhere}
%%   Although we have mentioned Mackey functors in connection with >>> it is not true that >>> is the category of $\MU_*\MU$-comodule valued Mackey functors. We suspect that if one studied $C_2$ equivariant motives, i.e. $\SH(B\mu_2)$, that a corresponding story would play out with $\MU_\R$ replaced by $\MU_{C_2}$. In this case the the modules over associated [tau] element would actually be comodule valued Mackey functors. 
%% \end{rmk}
%
%The proof of this theorem will occupy the next several pages.
%We begin with several lemmas whose proofs are short and abstract (Lemmas \ref{lem:gr-of-comparison}, \ref{lem:MU-koszul} and \ref{lem:uMU-koszul}).
%These lemmas puts us in a position where \Cref{lem:gr-R-bullet-computation} finishes the proof.
%The proof of \Cref{lem:gr-R-bullet-computation} is much more concrete, it proceeds via a direct examination of the object $R_\bullet$ constructed in [location].


\subsection{Twistings I: $t$-structures on filtered objects} \label{subsec:twists}\ 

%In the beginning of this section we highlighted the fact that the $(2n)^{\mathrm{th}}$ slice of $\MU_\R$ is a sum of copies of $\Sigma^{n \rho} \uZ$, which is \emph{not} in (a shift of) the heart of the usual $t$-structure. From this we were able to conclude that the obvious comparison functor between $C\ta$-modules and graded equivariant objects was inadiqute for our purposes.
%In this subsection, we discuss twisted $t$-structures on filtered and graded categories. Using this, we will produce ``twisting functors'' which will allow us to . We will take particular care to study how these twisting functors interact with symmetric monoidal structures.
%
We begin this subsection with a discussion of several natural $t$-structures on categories of graded and filtered objects, obtained by twisting a given $t$-structure by a Picard element.
In the filtered case our constructions are a straightforward generalization of a construction of Beilinson. Although these constructions are simple, they have the effect of modifying the symmetric monoidal structure on the heart of a category.
Next, we explicitly describe this modification at the level of the heart, where it admits a description in terms of an Euler characteristic. Read another way, there is an Euler characteristic obstruction to producing symmetric monoidal twisting functors. 
We close the subsection by producing symmetric monoidal twisting functors whenever this Euler characteristic obstruction vanishes.


\begin{dfn} \label{dfn:L-twist}
  For the remainder of this subsection, we fix the following data:
  a stable, presentably symmetric monoidal category $\mathcal{C}$
  equipped with a compatible $t$-structure\footnote{Here compatible means that a tensor product of objects which are $\geq 0$ is $\geq 0$ itself, and that the unit is $\geq 0$.}
  and a Picard element $\mathcal{L}$.

  We then define the $\mathcal{L}$-twisted $t$-structure on $\mathcal{C}^{\Gr}$ and $\mathcal{C}^{\Fil}$ as follows:\todo{Cite something to show this allowed?}
  \begin{itemize}
  \item An object $X_\bullet$ of $\mathcal{C}^{\Gr}$ is $\geq 0$ in the $\mathcal{L}$-twisted $t$-structure if, for all $n$, $\mathcal{L}^{\otimes -n} \otimes X_n$ is $\geq 0$ in the original $t$-structure.
  \item Similarly, an object $X_\bullet$ of $\mathcal{C}^{\Fil}$ is $\geq 0$ in the $\mathcal{L}$-twisted $t$-structure if $\mathcal{L} ^{\otimes -n} \otimes X_n$ is $\geq 0$ in the original $t$-structure, for all $n$.
  \end{itemize}
  We denote the hearts of these $t$-structures by $\CC^{\Gr, \mathcal{L}, \heartsuit}$ and $\CC^{\Fil, \mathcal{L}, \heartsuit}$ respectively.
  We denote the $i^{\mathrm{th}}$ $\mathcal{L}$-twisted $t$-structure homotopy objects by $\piu_{i} ^{\mathcal{L}, \Gr}$ and $\piu_{i} ^{\mathcal{L}, \Fil}$ respectively.
  When $\mathcal{L} = \o$, we frequently omit it from the notation.
\end{dfn}

We summarize the basic properties of this definition in the following lemma.

\begin{lem} \label{prop:twist-grade} \label{prop:twist-filt} \label{prop:twist-1-conn}
  In the graded case:
  \begin{enumerate}
  \item The $\mathcal{L}$-twisted $t$-structure on $\CC^{\Gr}$ is compatible with the symmetric monoidal structure.
  \item An object $X_\bullet$ is $< 0$ in the $\mathcal{L}$-twisted $t$-structure if and only if, for all $n$, $\mathcal{L}^{\otimes -n} \otimes X_n$ is $< 0$ in the original $t$-structure.
  \item The $t$-structure homotopy groups are determined by the following formula:
\[ (\underline{\pi}_0^{\mathcal{L}, \Gr} X_\bullet )_n \cong \mathcal{L}^{\otimes n} \underline{\pi}_0(\mathcal{L}^{\otimes -n} \otimes X_n). \]
  \end{enumerate}
  Assuming that $\mathcal{L} \geq 0$, we obtain similar results in the filtered case:
  \begin{enumerate}
  \item[(1')] The $\mathcal{L}$-twisted $t$-structure on $\CC^{\Fil}$ is compatible with the symmetric monoidal structure.
  \item[(2')] An object $X_\bullet$ is $< 0$ in the $\mathcal{L}$-twisted $t$-structure if and only if $\mathcal{L}^{\otimes -n} \otimes X_n < 0$.
  \item[(3')] The $t$-structure homotopy groups are given the by following formula:
  \[ (\underline{\pi}_0^{\mathcal{L}, \Fil} X_\bullet )_n \cong \mathcal{L}^{\otimes n} \underline{\pi}_0(\mathcal{L}^{\otimes -n} \otimes X_n). \]

  \end{enumerate}

  Under the stronger assumption that $\mathcal{L} \geq 1$, 
  the functor $\CC^{\Gr} \to \CC^{\Fil}$ which is right adjoint to the associated graded functor
  restricts to an equivalence of symmetric monoidal $1$-categories
  \[ \CC^{\Gr, \mathcal{L}, \heartsuit} \simeq \CC^{\Fil, \mathcal{L}, \heartsuit}. \]
\end{lem}

\begin{ntn}
  We write $- \otimes^{\heartsuit, \mathcal{L}} -$ for the tensor product induced on either $\CC^{\Gr, \mathcal{L}, \heartsuit}$ or $\CC^{\Fil, \mathcal{L}, \heartsuit}$. In the case that $\mathcal{L} = \o$, we simply write $- \otimes ^{\heartsuit} -$.
\end{ntn}

%The $\mathcal{L}$-twisted $t$-structure is easy to understand in the graded case:
%
%If we assume that $\mathcal{L} \geq 0$, then we may obtain similar results in the filtered case.

\begin{proof}
  The statements (1), (2) and (3) are all clear. We therefore restrct ourselves to the filtered case for the rest of the proof.

  From the expression of the tensor product as a Day convolution, we observe that a tensor product of connective objects has $n^{\mathrm{th}}$ term presented as a colimit over a diagram of connective objects tensored with $\mathcal{L}$ at least $n$ times. Since $\mathcal{L} \geq 0$ we may conclude that (1') holds.

  The expression for the homotopy objects in (3') follows from (2') in a straightforward way.
  We now prove (2').
  Suppose that $X_\bullet$ is an object such that $\mathcal{L}^{\otimes -n} \otimes X_n < 0$.
  We want to show that $X_\bullet < 0$, i.e. that it recieves no maps from $Y_\bullet \geq 0$.
  Using the condition that $\mathcal{L} \geq 0$, we learn that $[Y_n, X_m] = 0$ for all $n > m$, which is enough to imply that $[Y_\bullet, X_\bullet] = 0$.

  Suppose that $X_\bullet < 0$; then we would like to show that $\mathcal{L}^{\otimes -n} \otimes X_n < 0$.
  Associated to every $Y \geq 0$ in $\CC$ we can consider the filtered object $Y \otimes \mathcal{L}^{\otimes n}(n)$, which is depicted below:
  \[ \dots \to 0 \to 0 \to Y \otimes \mathcal{L} ^{\otimes n} \xrightarrow{\id} Y \otimes \mathcal{L} ^{\otimes n} \xrightarrow{\id} \dots \]
Here, the first nonzero object occurs at position $n$.
  Using the asssumption that $\mathcal{L} \geq 0$, we can conclude that $Y \otimes \mathcal{L}^{\otimes n}(n) \geq 0$. Now, on mapping out of this object we have
  \[ * \simeq \Omega^{\infty} \Map_{\CC^{\Fil}} ( Y \otimes \mathcal{L}^{\otimes n}(n), X) \simeq \Omega^\infty \Map_\CC ( Y \otimes \mathcal{L}^{\otimes n}, X_n), \]
  which implies that $\mathcal{L}^{\otimes -n} \otimes X_n < 0$ for all $n$, as desired.

  We now proceed to prove the final statement.
  Let $\tau : \o_{\CC^{\Fil}} (-1) \to \o_{\CC^{\Fil}}$ denote the shift map in $\CC^{\Fil}$. Then $C\tau$ admits an $\E_\infty$-algebra structure in $\CC^{\Fil}$ such that there is an equivalence of symmetric monoidal categories $\CC^{\Gr} \simeq \Mod(\CC^{\Fil}; C\tau)$ (cf. \Cref{prop:mod-fil}). Using the assumption that $\mathcal{L} \geq 0$, we learn that $\o_{\CC^{\Fil}}$ and $C\tau$ are $\geq 0$. As a consequence, we find that $\CC^{\Gr, \mathcal{L}, \heartsuit}$ may be identified with $\Mod(\CC^{\Fil, \mathcal{L}, \heartsuit}; \piu_0 ^{\mathcal{L}} C\tau)$.

  The proposition will therefore follow if we prove that $\piu_{0} ^{\mathcal{L}} \o_{\CC^{\Fil}} \to \piu_0 ^{\mathcal{L}} C\tau$ is an equivalence.
  By \Cref{prop:twist-filt}, we find that
  \[(\piu_{0} ^{\mathcal{L}, \Fil} \o_{\CC^{\Fil}})_n \cong \begin{cases} \mathcal{L}^{\otimes n} \pi_0 (\mathcal{L} ^{\otimes -n}) & \text{ if } n \leq 0, \\ 0 & \text{ otherwise.} \end{cases}\]
    Since $\mathcal{L} \geq 1$ by assumption, we find for $n < 0$ that $\mathcal{L}^{\otimes n} \pi_0 \mathcal{L} ^{\otimes -n} \cong 0$, whereas for $n=0$ we just get $\pi_0 \o_{\CC}$.

  On the other hand, we have that
  \[(\piu_0 ^{\mathcal{L}, \Fil} C\tau)_n \cong \begin{cases} \pi_0 \o_{\CC} & \text{ if } n = 0, \\ 0 & \text{ otherwise.} \end{cases}\]
  It follows from this that the map $\piu_{0} ^{\mathcal{L}, \Fil} \o_{\CC^{\Fil}} \to \piu_0 ^{\mathcal{L}, \Fil} C\tau$ is an equivalence, as desired.
\end{proof}

We can now profitably define the twist functors.

\begin{cnstr} \label{cnstr:twist}
  Given a symmetric monoidal Grothendieck abelian $1$-category $\mathcal{A}$ and an invertible object $a \in \mathcal{A}$, we can construct a monoidal functor
  $i_a : \Z \to \mathcal{A}$ which sends $1$ to $a$.\footnote{By symmetric monoidal Grothendieck abelian $1$-category, we mean a commutative algebra object of the symmetric monoidal category $\mathrm{Groth}_1$ discussed in \Cref{rec:groth-ab}. In other words, we assume that the tensor product commutes with colimits in each variable.}
  If the swap map $s_{a,a}$ is the identity, then we can make this functor symmetric monoidal.
  Tensoring up with $\Abh$, we obtain a monoidal (or symmetric monoidal) functor
  \[ i_a : \Ab^{\Gr, \heartsuit} \xrightarrow{\o(1) \mapsto a} \mathcal{A}. \]
  
  In our situation of interest, we take $\mathcal{A} = \CC^{\Gr, \mathcal{L}, \heartsuit}$ and $a = (\piu_0 \o) \otimes \mathcal{L}(1)$. Then we can define the functor $\twist^{\mathcal{L}}$ as the composite
  \[ \CC^{\Gr, \heartsuit} \cong \CC^\heartsuit \otimes \Ab^{\Gr, \heartsuit} \xrightarrow{\mathrm{Id} \otimes i_{(\piu_0 \o) \otimes \mathcal{L}(1)}} \CC^\heartsuit \otimes \CC^{\Gr, \mathcal{L}, \heartsuit} \xrightarrow{\otimes} \CC^{\Gr, \mathcal{L}, \heartsuit}. \]
  On objects this has the effect of sending $\{X_n\}$ to $\{X_n \otimes \mathcal{L}^{\otimes n}\}$.
  It is easy to see that this is a monoidal equivalence between $\CC^{\Gr, \heartsuit}$ and $\CC^{\Gr, \mathcal{L}, \heartsuit}$, and is further symmetric monoidal if $s_{a, a}$ is the identity. 
\end{cnstr}

The above construction succinctly explains how much the symmetric monoidal structure on $\CC^{\Gr, \mathcal{L}, \heartsuit}$ has been twisted in terms of the quantity $s_{a, a}$. We now give a description of this quantity in terms of the following standard definition:

\begin{dfn}
  Given a dualizable object $X$,
  with dual $X^\vee$, unit $\eta$, and counit $\epsilon$,
  the trace of an endomorphism $f: X \to X$ is the element $\mathrm{tr}(f) \in \pi_0 \End(\o)$ given by the composite
  \[ \o \xrightarrow{\eta} X \otimes X^\vee \xrightarrow{ f \otimes X^\vee} X \otimes X^\vee \xrightarrow{s_{X,X^\vee}} X^{\vee} \otimes X \xrightarrow{\epsilon} \o. \]
  The trace of the identity is called the Euler characteristic and denoted $\chi(X)$.
\end{dfn}

A simple diagram chase tells us that the Euler characteristic is multiplicative and defines an $\End(\o)$-linear map $\mathrm{tr} : \End(X) \to \End(\o)$. Moreover, one can compute that the trace of the swap map $s_{X,X}$ is equal to $\chi(X)$. In the case where $X$ is invertible the trace map is an isomorphism, and we can use this to conclude that the condition $s_{a,a} = \mathrm{Id}_{a \otimes a}$ in \Cref{cnstr:twist} is equivalent to asking that $\chi(a) = 1$. Specializing further to the case $a = (\piu_0 \o) \otimes \mathcal{L}(1)$, we note that the diagram computing the Euler characteristic of $(\piu_0 \o) \otimes \mathcal{L}(1)$ is $\piu_0^{\mathcal{L},\Gr}$ applied to the diagram computing the Euler characteristic of $\mathcal{L}(1)$ in the ambient category $\CC^{\Gr}$. This implies that
\[ \chi(a) = \chi(\mathcal{L}(1)) = \chi(\mathcal{L})\chi(\o(1)) = \chi(\mathcal{L}) \cdot 1. \]
Thus, we have proved:

\begin{lem} \label{prop:twist}
  The functor $\twist^{\mathcal{L}}$ can be made symmetric monoidal if $\chi(\mathcal{L}) = 1$.
\end{lem}

%% \begin{lem}
%%   The assignment $\{X_n\} \mapsto \{X_n \otimes \mathcal{L}^{\otimes n}\}$ determines an automorphism of $\twist^{\mathcal{L}}$ of $\CC^{\Gr}$ which is $t$-exact if we equip the source with the standard $t$-structure and the target with the $\mathcal{L}$-twisted $t$-structure. 
%% \end{lem}

%% \begin{proof} Clear from \Cref{dfn:L-twist} and \Cref{prop:twist-grade}. \end{proof}

%% The question we would like to answer is ``when can these functors be equipped with a symmetric monoidal structure?''. The first obstruction to this comes from the Euler characteristic, which we define below.

Conversely, if we assume that $\twist^{\mathcal{L}}$ is symmetric monoidal, then we can make the following Euler characteristic computation:
\[ 1 = \twist^{\mathcal{L}} (1) = \twist^{\mathcal{L}}( \chi (\o(1))) = \chi ( \twist^{\mathcal{L}}(\o(1))) = \chi ( \mathcal{L}(1) ) = \chi (\mathcal{L}). \]
%% In fact, we can show that at the level of the heart this is the only obstruction to making the functor $\twist^{\mathcal{L}}$ symmetric monoidal. 

%% \begin{lem} \label{prop:twist}
%%   Since $\twist^{\mathcal{L}}$ is $t$-exact it naturally restricts to a functor,
%%   \[ \twist^{\mathcal{L}} : \CC^{\Gr, \heartsuit} \simeq \CC^{\Gr, \mathcal{L}, \heartsuit}. \]
%%   This functor may naturally be given the structure of a monoidal functor.
%%   Moreover, it may be given a symmetric monoidal structure if $\chi(\mathcal{L}) = 1$.
%% \end{lem}

%% \begin{proof}
%%    To give $\twist^{\mathcal{L}}$ the structure of a monoidal functor, we must specify natural isomorphisms
%%    \[\{X_n \otimes \mathcal{L}^{\otimes n}\} \otimes^{\heartsuit, \mathcal{L}} \{Y_n \otimes \mathcal{L}^{\otimes n}\} \cong \{(X_\bullet \otimes^{\heartsuit} Y_\bullet)_n \otimes \mathcal{L}^{\otimes n}\},\]
%%    i.e. specify natural isomorphisms
%%    \[\piu_0 ^{\mathcal{L}} (\{X_n \otimes \mathcal{L}^{\otimes n}\} \otimes \{Y_n \otimes \mathcal{L}^{\otimes n}\}) \cong \{(\piu_{0} (X_\bullet \otimes Y_\bullet))_n \otimes \mathcal{L}^{\otimes n}\}.\]
%%    Now, by \Cref{prop:twist-grade}, we have a natural isomorphism
%%    \[\{(\piu_{0} (X_\bullet \otimes Y_\bullet))_n \otimes \mathcal{L}^{\otimes n}\} \cong \piu_0 ^{\mathcal{L}} (\{(X_\bullet \otimes Y_\bullet)_n \otimes \mathcal{L}^{\otimes n}\}).\]
%%    It therefore suffices to specify natural isomorphisms
%%    \[\{X_n \otimes \mathcal{L}^{\otimes n}\} \otimes \{Y_n \otimes \mathcal{L}^{\otimes n}\} \cong \{(X_\bullet \otimes Y_\bullet)_n \otimes \mathcal{L}^{\otimes n}\}\]
%%    in the homotopy $1$-category $h\CC$ of $\CC$. From this point on we work in the symmetric monoidal $1$-category $h\CC$.

%%    Expanding out the definitions, we find that we need to provide natural isomorphisms
%%    \[\{\bigoplus_{a+b=n} (X_a \otimes \mathcal{L}^{\otimes a}) \otimes (Y_b \otimes \mathcal{L}^{\otimes b})\} \cong \{(\bigoplus_{a+b=n} X_a \otimes Y_b) \otimes \mathcal{L}^{\otimes n}\}.\]
%%    We now provide, for each $a$ and $b$, an equivalence
%%    \[(X_a \otimes \mathcal{L}^{\otimes a}) \otimes (Y_b \otimes \mathcal{L}^{\otimes b}) \cong (X_a \otimes Y_b) \otimes \mathcal{L}^{\otimes a+b},\]
%%    namely the swap map $s_{\mathcal{L}^{\otimes a}, Y_b}$ (here and throughout the rest of the proof we will suppress the associativity isomorphisms).

%%    It is not hard to check that this equips $\twist^{\mathcal{L}}$ with the structure of a monoidal functor.
%%    To check that it is in fact a symmetric monoidal functor, we need to check that the following diagram commutes for each $a$ and $b$:
%%    \begin{center}
%%      \begin{tikzcd} [column sep = 1.5in, row sep = large]
%%        (X_a \otimes \mathcal{L} ^{\otimes a}) \otimes (Y_b \otimes \mathcal{L}^{\otimes b}) \ar[r,"s_{(X_a \otimes \mathcal{L} ^{\otimes a}), (Y_b \otimes \mathcal{L}^{\otimes b})}"] \ar[d,"s_{\mathcal{L}^{\otimes a}, Y_b}"] & (Y_b \otimes \mathcal{L} ^{\otimes b}) \otimes (X_a \otimes \mathcal{L}^{\otimes a}) \ar[d, "s_{\mathcal{L}^{\otimes b}, X_a}"] \\
%%        (X_a \otimes Y_b) \otimes \mathcal{L}^{\otimes a+b} \ar[r,"s_{X_a, Y_b}"] & (Y_b \otimes X_a) \otimes \mathcal{L}^{\otimes a+b}.
%%      \end{tikzcd}
%%    \end{center}
%%    Tracing around the the bottom of the diagram, we obtain the map
%%    \[s_{X_a, \mathcal{L}^{\otimes b}} \circ s_{X_a, Y_b} \circ s_{\mathcal{L}^{\otimes a}, Y_b}.\]
%%    On the other hand, going across the top, we have
%%    \begin{align*}
%%      &s_{(X_a \otimes \mathcal{L} ^{\otimes a}), (Y_b \otimes \mathcal{L}^{\otimes b})} \\
%%      &= s_{(X_a \otimes \mathcal{L} ^{\otimes a}), \mathcal{L}^{\otimes b}} \circ s_{(X_a \otimes \mathcal{L} ^{\otimes a}), Y_b} \\
%%      &=  s_{X_a, \mathcal{L}^{\otimes b}} \circ s_{\mathcal{L} ^{\otimes a}, \mathcal{L}^{\otimes b}} \circ s_{X_a, Y_b} \circ s_{\mathcal{L} ^{\otimes a}, Y_b}.
%%    \end{align*}
%%    Now, we have $s_{\mathcal{L}^{\otimes a}, \mathcal{L}^{\otimes b}} = (s_{\mathcal{L}, \mathcal{L}})^{\circ (ab)}$ and therefore is will suffice to show that $s_{\mathcal{L}, \mathcal{L}}$ is the identity.

%%    The trace defines a $\pi_0\o$-linear isomorphism between $\pi_0End(\mathcal{L}^{\otimes 2})$ and $\pi_0\o$ which sends the identity to $\chi(\mathcal{L}^{\otimes 2}) = \chi( \mathcal{L} )^2$. One can then directly compute that the trace of $s_{\mathcal{L}, \mathcal{L}}$ is equal to $\chi ( \mathcal{L} )$. Now, using the assumption that $\chi (\mathcal{L}) = 1$ we may conclude.
%% \end{proof}

\begin{rmk}
  Examining the failure of $\twist^{\mathcal{L}}$ to be symmetric monoidal, we find that in general the symmetric monoidal structure on $\CC^{\Gr, \mathcal{L}, \heartsuit}$ is twisted by the Euler characteristic $\chi(\mathcal{L})$.
\end{rmk}
%
%\begin{rec}
%  Given a presentably symmetric monoidal category $\CC$, we may consider the Picard category $\Pic(\CC)$. This is the full subcategory
%
%  If $\CC$ is a $1$-category, then $\pi_n \pic (\CC) = 0$ for $n \neq 0,1$. Moreover, $\pi_0 \pic(\CC)$ is the Picard group of $\CC$ and $\pi_1 \Pic(\CC) \cong \aut(\o)$.
%\end{rec}
%
%\begin{prop}
%  The cofiber of the map $\pic (\CC^{\heartsuit}) \to \pic (\CC^{\Gr, \mathcal{L}, \heartsuit})$ is equivalent to $\Z$.
%\end{prop}
%
%\begin{proof}
%  Since $\CC^{\heartsuit} \to \CC^{\Gr, \mathcal{L}, \heartsuit}$ is fully faithful, the induced map  $\pic (\CC^{\heartsuit}) \to \pic (\CC^{\Gr, \mathcal{L}, \heartsuit})$
%\end{proof}
%
%\begin{prop}
%  The natural symmetric monoidal functor $\mathcal{P} (\Pic (\CC^{\Gr, \mathcal{L}, \heartsuit})) \otimes_{\mathcal{P} (\Pic (\CC^{\heartsuit}))} \CC^{\heartsuit} \to \CC^{\Gr, \mathcal{L}, \heartsuit}$ is an equivalence.
%\end{prop}
%
%As a consequence, the symmetric monoidal structure on $\CC^{\Gr, \mathcal{L}, \heartsuit}$ only depends on the map of spectra $\pic (\CC^{\heartsuit}) \to \pic (\CC^{\Gr, \mathcal{L}, \heartsuit})$.
%
%It follows that this map is determined by a map $\Z \to \Sigma \pic(\CC^{\heartsuit})$.
%\begin{lem}
%  Let $A$ and $B$ denote abelian groups, viewed as Eilenberg-Maclane spectra. Let $B[2] \subset B$ denote the subgroup of $2$-torsion elements. Then there are isomorphisms:
%  \[ [A,\Sigma B] \cong \Ext_{\Z} (A,B) \]
%  \[ [A, \Sigma^2 B] \cong \Hom_{\Z} (A,B[2]). \]
%\end{lem}
%
%We therefore need to understand the structure of $\pic(\CC^{\heartsuit})$.
%This has two nonzero homotopy groups: $\pi_0 \pic(\CC^{\heartsuit})$ is the Picard group of $\CC^{\heartsuit}$, and $\pi_1 \pic(\CC^{\heartsuit}) \cong \aut(\o_{\CC^{\heartsuit}})$.
%The single nonzero $k$-invariant $\pi_0 \pic(\CC^{\heartsuit}) \to \Sigma^{2} \pi_1 \pic(\CC^{\heartsuit}) \simeq \Sigma^{2} \aut(\o_{\CC^{\heartsuit}})$
%corresponds to the trace map $\Tr : \pi_0 \pic(\CC^{\heartsuit}) \to \aut(\o_{\CC^{\heartsuit}})[2]$ sending a Picard element $\mathcal{L}$ to its trace $\Tr(\mathcal{L})$.
%
%
%  The negative objects are easy to recognize in the graded case.
%  If we assume that $\mathcal{L}$ is $\geq 0$, then the negative are similarly easy to determine,
%  \begin{itemize}
%  \item An object $X_\bullet$ of $\mathcal{C}^{\Gr}$ is $\leq 0$ in the $\mathcal{L}$-twisted $t$-structure if and only if for all $n$, $\mathcal{L}^{\otimes -n} \otimes X_n$ is $\leq 0$ in the original $t$-structure.
%  \item Assuming that $\mathcal{L}$ is $\geq 0$, an object $X_\bullet$ of $\mathcal{C}^{\Fil}$ is $\leq 0$ in the $\mathcal{L}$-twisted $t$-structure if and only if its underlying graded object satisfies the same.
%  \end{itemize}
%
%In the graded case we have a simple formula for homotopy group functor,
%\[ (\underline{\pi}_0^{\mathcal{L}, \Gr} (X_\bullet) )_n \cong \mathcal{L}^{\otimes n} \underline{\pi}_0(\mathcal{L}^{\otimes -n} \otimes X_n) \]
%In the filtered case we have a similar formula as long as $\mathcal{L}$ is $\geq 0$,
%\[ (\underline{\pi}_0^{\mathcal{L}, \Fil} (X_\bullet) )_n \cong \mathcal{L}^{\otimes n} \underline{\pi}_0(\mathcal{L}^{\otimes -n} \otimes X_n) \]
%If we make the stronger assumption that $\mathcal{L}$ is $\geq 1$, then the shift map on an object of the heart becomes zero (up to contractible ambiguity) so $\underline{\pi}_0^{\mathcal{L}, \Fil}$ factors through the graded version of the heart with an equivalence
%\[ \underline{\pi}_0^{\mathcal{L},\Gr}(\Gr(-)) \cong \underline{\pi}_0^{\mathcal{L},\Fil}(-). \]
%Altogether, this lets us conclude that as long as $\mathcal{L}$ is $\geq 1$, there are equivalences
%\[ \mathcal{C}^{\Fil, \mathcal{L}\heartsuit} \cong \mathcal{C}^{\Gr, \mathcal{L}\heartsuit} \cong (\mathcal{C}^\heartsuit)^{\Gr}. \]
%where the second equivalence sends $X_\bullet$ to $\mathcal{L}^{\otimes \bullet} \otimes X_\bullet$.
%
%In the situation of \cite{dfn:L-twist} the compatibility of the $t$-structure with the tensor product turns $\mathcal{C}^\heartsuit$ into a symmetric monoidal $1$-category. Similarly, the assumption that $\mathcal{L}$ is $\geq 0$ ensures that the $\mathcal{L}$-twisted $t$-structure is compatible with the tensor product on filtered/graded objects. Using the equivalence of 1-categories from the previous paragraph, we have constructed a symmetric moniodal structure on $(\mathcal{C}^\heartsuit)^{\Gr}$ for each suitable $\mathcal{L}$. For later use we would like to study these symmetric monoidal structures.
%
%A monoidal $\mathcal{C}^\heartsuit$-linear functor with source $\mathcal{C}^{\heartsuit,\Gr}$ is equivalent to the choice of Picard element in the target (to which $\o(1)$ is sent). From this, we can conclude that the $\mathcal{L}$-twisted symmetric monoidal structure on $\mathcal{C}^{\heartsuit, \Gr}$ has the same underlying monoidal structure as the untwisted case.
%
%\begin{lem}
%  If we consider $\mathcal{C}^{\heartsuit, \Gr}$ as a monoidal $\mathcal{C}^\heartsuit$-algebra, then symmetric monoidal refinements of the monoidal structure are in bijective correspondence with 2-torsion elements of $\pi_0(\o)^\times$. The correspondence sends a symmetric monoidal refinement to $Tr(\o(1))$.
%\end{lem}
%
%\begin{proof}
%  \todo{write this.}
%\end{proof}
%

\begin{exm}
  If we take $\mathcal{C} = \Sp$ with its usual $t$-structure, then the $\Ss^1$-twisted category $\Sp^{\Gr, \Ss^1, \heartsuit}$ is equivalent to the symmetric monoidal category of graded abelian groups with the Koszul sign convention, since $\chi(\Ss^1) = -1$.

  On the other hand, $\Sp^{\Gr, \Ss^2, \heartsuit}$ is symmetric monoidally equivalent to the category of graded abelian groups by \Cref{prop:twist}, since $\chi(\Ss^2) = 1$. \todo{Define $\chi$ above, note that it is $2$-torsion for a Picard element.}
\end{exm}
%
%In [Ctau], the fact that $S^2$-twisted graded abelian groups are equivalent to graded abelian groups is implicitly used in proving their symmetric monoidal equivalence.
%\todo{Should we throw some shade about lazy non-proofs here?}

%
%Extending the example above, the case of interest for this paper is $\Ss^{\rho}$-twisted graded $\uZ$-module Mackey functors. Observe that since $[C_2] - 1$ maps to $1$ in $\uZ$ the twist is trivial (like the case over $\C$).

\begin{lem} \label{lem:D-twist}
  Suppose that $\CC$ is equivalent to the derived category of its heart $\CC^{\heartsuit}$ as a symmetric monoidal category with compatible $t$-structure. Then there exists an equivalence of monoidal categories $\twist^{\mathcal{L}} : \CC^{\Gr} \simeq \CC^{\Gr}$ making the following diagram of monoidal functors commute:
  \begin{center}
    \begin{tikzcd}
      \CC^{\Gr, \heartsuit} \ar[r, "\twist^{\mathcal{L}}"] \ar[d] & \CC^{\Gr, \mathcal{L}, \heartsuit} \ar[d] \\
      \CC^{\Gr} \ar[r, "\twist^{\mathcal{L}}"] & \CC^{\Gr}.
    \end{tikzcd}
  \end{center}
  If $\chi(\mathcal{L}) = 1$, then the above square naturally lifts to a diagram of symmetric monoidal functors.
  On objects, $\twist^{\mathcal{L}} : \CC^{\Gr} \to \CC^{\Gr}$ is given by $\{X_n\} \mapsto \{X_n \otimes \mathcal{L}^{\otimes n}\}$.
\end{lem}

\begin{proof}
  The assumption implies that $\CC^{\Gr}$, equipped with the usual $t$-structure, is equivalent as a symmetric monoidal category with compatible $t$-struture to $\mathcal{D}(\CC^{\Gr, \heartsuit})$.
  It also implies that $\CC^{\Gr}$, when equipped with the $\mathcal{L}$-twisted $t$-structure, is equivalent as a symmteric monoidal category with compatible $t$-struture to $\mathcal{D}(\CC^{\Gr, \mathcal{L}, \heartsuit})$.

  It follows that the equivalence of monoidal $1$-categories $\twist^{\mathcal{L}} : \CC^{\Gr, \heartsuit} \to \CC^{\Gr, \mathcal{L}, \heartsuit}$ of \Cref{prop:twist} determines a compatible equivalence of monoidal categories $\twist^{\mathcal{L}}: \CC^{\Gr} \to \CC^{\Gr}$, via taking derived categories.
  If $\chi(\mathcal{L}) = 1$, these are equivalences of symmetric monoidal categories by \Cref{prop:twist}.
\end{proof}



\subsection{Twistings II: the slice filtration} \label{subsec:slice}\ 

In this short subsection, we specialize the material from the previous subsection to the case of interest. This means we look at $\Sp_{C_2, i2}$ with Picard element $\Ss_2 ^\rho$.
In order to connect this with $\MU_{\R,2}$ and the material from \Cref{sec:top}, we relate the $\rho$-twisted $t$-structure to the regular slice filtration.

\begin{exm}
  If we take $\mathcal{C}= \Sp_{C_2}$, then the following Euler characteristic computations provide the three possible nontrivial twists on graded Mackey functors:
  \[ \chi(\Ss^{\sigma}) = 1 - [C_2], \qquad \chi(\Ss^1) = -1, \qquad \chi(\Ss^{\rho}) = [C_2] - 1. \]
  To verify this, one uses the compatibility of the Euler characteristic with the symmetric monoidal functors $\Phi^{e}$ and $\Phi^{C_2}$.
  %Each of these computations can be accomplished using the compatibility of Euler characteristic with symmetric monoidal functors applied to $\Phi^e$ and $\Phi^{C_2}$.
\end{exm}

\begin{exm}
  Write $\uAb_{i2}^{\Gr, \rho, \heartsuit}$ for $\uAb_{i2}^{\Gr, \Sigma^{\rho} \uZt, \heartsuit}$.
  Then, since the image of $\chi(\Ss^{\rho})$ in $\pi_0 \uZt$ is $1$,
  there is a symmetric monoidal equivalence
  \[\twist^{\rho} : \uAb_{i2}^{\Gr, \heartsuit} \xrightarrow{\simeq} \uAb_{i2}^{\Gr, \rho, \heartsuit},\]
  by \Cref{prop:twist}. Applying \Cref{lem:D-twist}, we can then upgrade this to a diagram of symmetric monoidal functors
  \begin{center}
    \begin{tikzcd}
      \uAb_{i2} ^{\Gr, \heartsuit} \ar[r, "\twist^{\rho}"] \ar[d] & \uAb_{i2} ^{\Gr, \rho, \heartsuit} \ar[d] \\
      \uAb_{i2} ^{\Gr} \ar[r, "\twist^{\rho}"] & \uAb_{i2} ^{\Gr}.
    \end{tikzcd}
  \end{center}
\end{exm}

\todo{We're gonna need to either explain more about how the slices work, cite somebody or write more shit here.}
\todo{We need to fix a convention does pi n live in deg n (composite of two truncation functors) or in deg 0.}

\begin{cnv}
  From this point on, our category $\CC$ will be one of $\Sp_{C_2, i2}$, $M(C_2)_{i2}$, $\uAb_{i2}$ or $\Ab_{i2}$ with its usual $t$-structure, and the Picard element we work with will be $\Ss^{\rho} _2, \Sigma^{\rho} \piu_{0} ^{C_2} \Ss_2, \Sigma^{\rho} \uZt$ or $\Sigma^{2} \Zt$, respectively. In order to reduce the notational burden, we denote the former three Picard objects by $\rho$ and the final Picard object by $2$ in superscripts. So, in the example of $\Sp_{C_2, i2}$, we will write $\tau_{\geq 0} ^\rho$ for the connective cover, $\piu_0 ^\rho$ for the $0^{\mathrm{th}}$ homotopy object and $\Sp_{C_2, i2} ^{\Fil, \rho, \heartsuit}$ for the heart of $\Sp_{C_2, i2} ^{\Fil}$.
%% We equip the category $\Sp_{C_2, i2} ^{\Fil}$ with the $\Ss^{\rho} _{2}$-twisted $t$-structure, and let 
%%   We denote the heart by $\Sp_{C_2, i2} ^{\Fil, \rho, \heartsuit}$, and note that this is naturally equivalent to $\Sp_{C_2, i2} ^{\Gr, \rho, \heartsuit}$ by \Cref{prop:twist-1-conn}.
%%   This is further equivalent to $M(C_2)^{\Gr, \rho, \heartsuit}$, since $M(C_2)$ is equivalent to the category of modules over $\piu_0 ^{C_2} \Ss_2$ in $\Sp_{C_2, i2}$.
\end{cnv}

This key result of this section is the following:

\begin{lem}\label{prop:twist-slice}
  Let $E \in \Sp_{C_2, i2}$. Then
  \[\tau_{\geq 0} ^{\rho} Y(E) \simeq \left( \dots \to P_4 E \to P_2 E \to P_0 E \to P_{-2} E \to P_{-4} E \to \dots \right), \]
  where $Y$ is the functor that sends an object to the associated constant filtered object. 
  If we further suppose that $\piu_{n\rho-1} ^{C_2} E = 0$ for all $n$, there is a natural equivalence
  \[\Gr (\tau_{\geq 0} ^{\rho} Y(E)) \simeq \piu_0 ^{\rho} Y(E) \cong \{\Sigma^{n\rho} \piu_{n\rho} ^{C_2} E\}.\]
  Finally, if $\piu_{n\rho} E$ is a constant Mackey functor for all $n$, then
  $ \piu_0 ^{\rho} Y(E) $
  factors through the full subcategory
  \[\uAb_{i2} ^{\Gr, \rho, \heartsuit} \subset M(C_2)_{i2} ^{\Gr, \rho, \heartsuit} \simeq \Sp_{C_2, i2} ^{\Fil, \rho, \heartsuit}.\]
\end{lem}

\begin{proof}
  This lemma has three statements.
  For the first statement, it suffices to note the following:
  \begin{enumerate}
    \item A $C_2$-spectrum is $t$-structure $0$-connective if and only if it is regular slice $0$-connective.
\item A $C_2$-spectrum $E$ is regular slice $2n$-connective if and only if $E \otimes \Ss^{-n\rho}$ is regular slice $0$-connective.
  \end{enumerate}

  For the second statement, since $\Ss^{\rho} _2 \geq 1$, we can use the final statement of \Cref{prop:twist-1-conn} to obtain a canonical map
  \[\Gr (\tau_{\geq 0} ^{\rho} Y(E)) \to \piu_0 ^{\rho} Y(E).\]
  On the one hand, we have
  $ \piu_0 ^\rho Y(E) \cong \{\Sigma^{n\rho} \piu_{n\rho} ^{C_2} E\} $
  by \Cref{prop:twist-grade}(3'). \todo{I dont 100\% follow this argument.}
  Since by assumption the odd slices of $E$ vanish, we learn that
  the $(2n)^{\mathrm{th}}$ slice of $E$ is equivalent to $\Sigma^{n\rho} \piu_{n\rho} ^{C_2} E$ and
  that there are cofiber sequences
  \[ P_{2n+2} E \to P_{2n} E \to \Sigma^{n\rho} \piu_{n\rho} ^{C_2} E. \]
  The desired equivalence now follows from the first statement.

  The third statement now follows from the fact that any discrete $C_2$-Mackey functor which is constant admits a (necessarily unique) $\uZ$-module structure.
\end{proof}




\subsection{Reductions} \label{subsec:reds}\ 

In this subsection, we express each side of the equivalence (\ref{eqn:cta}) of \Cref{thm:Cta} as the category of modules over some commutative algebra in $\Sp_{C_2, i2} ^{\Fil}$.
Almost everything in this subsection is a standard application of Lurie's theory of higher algebra.
We have extracted the necessary material in \Cref{app:boilerplate}, where we present it in a digested form. 

We begin with the left-hand-side of (\ref{eqn:cta}).

\begin{lem} \label{lem:gr-of-comparison}
  The cofiber of $\ta : \Ss_2^{0,0,-1} \to \Ss_2^{0,0,0}$ is a commutative algebra in $\SHR_{i2}$. Furthermore, there is an equivalence of symmetric monoidal categories
  \[ \mathrm{Mod}(\SHRat; C\ta) \simeq \mathrm{Mod}( \Sp_{C_2, i2}^{\Gr}; \Gr( R_\bullet ) ), \]
	where $R_{\bullet} =\Tot^{*}\left( P_{2\bullet} \MU_{\mathbb{R},2}^{\otimes *+1} \right)$ 
\end{lem}

\begin{proof}
  Recall that the shift map $\Ss_2(-1) \to \Ss_2$ in $\Sp_{C_2, i2}^{\Fil}$ is denoted by $\tau$.
    Tensoring the equivalence of \Cref{exm:fil-ctau} with $\Sp_{C_2, i2}$ and applying \Cref{lem:tensor-mod}, we see that there is an equivalence of presentably symmetric monoidal categories
    \[ \Sp_{C_2, i2} ^{\Gr} \simeq \Mod ( \Sp_{C_2, i2} ^{\Fil} ; C\tau ). \]

    Under the symmetric monoidal functor $i^* : \Sp_{C_2, i2} ^{\Fil} \to \SHRat$ constructed in \Cref{prop:top-diagram}, the commutative algebra $C\tau$ maps to $C\ta$. Therefore, using \Cref{lem:tensor-mod}, we have 
    \[ \mathrm{Mod}(\SHRat; C\ta) \simeq \Sp_{C_2, i2}^{\Gr} \otimes_{\Sp_{C_2, i2} ^{\Fil}} \SHRat. \]
    Using \Cref{thm:filt-model}, there is an equivalence of presentably symmetric monoidal categories under $\Sp_{C_2, i2}^{\Fil}$ between $\SHRat$ and $\mathrm{Mod}( \Sp^{\Fil}_{C_2,i2} ; R_\bullet )$. Tensoring down along the associated graded ring map and using \Cref{lem:tensor-mod}, we obtain equivalences
    \begin{align*}
      \Sp_{C_2, i2}^{\Gr} \otimes_{\Sp_{C_2, i2}^{\Fil}} \SHRat
      %&\cong \Sp^{\Gr} \otimes_{\Sp^{\Fil}} \mathrm{Mod}( \Sp_{C_2,i2}^{\Fil} ; \Gamma_* \Ss_2 ) \\
      &\simeq \Sp_{C_2,i2}^{\Gr} \otimes_{\Sp_{C_2,i2}^{\Fil}} \mathrm{Mod}( \Sp_{C_2,i2}^{\Fil} ; R_\bullet ) \\
      &\simeq \mathrm{Mod}( \Sp_{C_2,i2}^{\Fil} ; \Gr(R_{\bullet}) ). \qedhere
    \end{align*}
      %
  %% Note that the shift map $\tau: \o(1) \to \o$ in $\Fil(\Sp_{C_2, i2})$ goes to $\ta$ under the colimit-preserving symmetric monoidal functor $c_{\C/\R} : \Fil(\Sp_{C_2, i2}) \to \SH(\R)^{\at} _{i2}$. Since $C\tau$ is an $\E_{\infty}$-ring in $\Fil(\Sp_{C_2, 2})$, it follows that its image $C\ta$ under $c_{\C/\R}$ is as well.

  
  %% %% \[ \mathrm{Mod}(\SHgc_{\mathrm{i}p}; C\ta) \simeq \Gr(\Sp) \otimes_{\Fil(\Sp)} \SHgc_{\mathrm{i}p} \simeq \Gr(\Sp) \otimes_{\Fil(\Sp)} \mathrm{Mod}(\Fil(\Sp_{C_2}); R_\bullet) \simeq \mathrm{Mod}(\Gr(\Sp_{C_2}); R_\bullet) \]
  %%   \[ \Gr(\Sp_{C_2, i2}) \otimes_{\Fil(\Sp_{C_2, i2})} \SH(\R)^{\at}_{i2} \simeq \Gr(\Sp_{C_2, i2}) \otimes_{\Fil(\Sp_{C_2, i2})} \mathrm{Mod}(\Fil(\Sp_{C_2,2}); \Gamma_* (\Ss_2)) \]
  %%   symmetric monoidal categories. Using \Cref{lem:tensor-mod}, we can simplify each side of this equivalence to obtain:
  %%   \[ \Gr(\Sp_{C_2, i2}) \otimes_{\Fil(\Sp_{C_2, i2})} \SH(\R)^{\at}_{i2} \simeq \mathrm{Mod}(\SH(\R)^{\at}_{i2}; C\ta), \]
  %%   \[ \Gr(\Sp_{C_2, i2}) \otimes_{\Fil(\Sp_{C_2, i2})} \mathrm{Mod}(\Fil(\Sp_{C_2}); (\Gamma_* (\Ss_2)) \simeq \mathrm{Mod}(\Gr(\Sp_{C_2}); \Gr(\Gamma_* (\Ss_2))), \]
  %%   where in the first case we have used the fact that $\Mod(\Fil(\Sp_{C_2, i2}); C\tau) \simeq \Gr(\Sp_{C_2,i2})$.
\end{proof}

Now we proceed to the right-hand-side of (\ref{eqn:cta}).

\begin{lem} \label{lem:omega-euler-one}
  The Euler characteristic of $\omega_{\G/\Mfg} \in \IndCoh(\Mfg)^{\heartsuit}$ is equal to 1.
  %% one: $\chi(\omega_{\G/\Mfg}) = 1$.
  As a consequence, \Cref{cnstr:twist} provides a symmetric monoidal left adjoint
  \[ p^* : \Ab^{\Gr, \heartsuit} \to \IndCoh(\Mfg)^{\heartsuit} \]
  %% \subset \IndCoh(\Mfg)
  which sends $\Z(1)$ to $\omega_{\G/\Mfg}$.
\end{lem}

\begin{proof}
  Let $L$ denote the Lazard ring. Then the flat cover $\Spec(L) \to \Mfg$ determines a pullback map $\IndCoh(\Mfg) \to \Mod_{L}$. This determines an injective pullback map
  $\pi_0 \End(\mathcal{O}_{\Mfg}) \to \pi_0 \End(L)$.
  Since the pullback of $\omega_{\G/\Mfg}$ is equivalent to $L$ we learn that $\chi(\omega_{\G/\Mfg}) = 1$.
  %% The latter group is isomorphic to $\Z ^\times$, which the former group contains, so the map is an equivalence.
  %% It follows that $\chi(\omega_{\G/\Mfg})$ is equal to the Euler characteristic of the pullback of $\omega_{\G/\Mfg}$ to $\Mod_{L}$. But this pullback is just $L$ itself, which implies that its Euler characteristic must be $1$, as desired.
  %% For the second statement, it is clear that one may construct a monoidal functor $\Z \to \IndCoh(\Mfg)^{\heartsuit}$ sending $1$ to $\omega_{\G/\Mfg}$. To check that it is symmetric monoidal, one has to check that it sends the swap automorphism of $2 = 1+1$, which is the identity, to the swap isomorphism of $\omega_{\G/\Mfg} ^2 \simeq \omega_{\G/\Mfg} \otimes \omega_{\G/\Mfg}$. This is precisely the condition that $\chi(\omega_{\G/\Mfg}) = 1$.
\end{proof}

Since $p^*$ sends a family of compact dualizable objects (the powers of $\Z(1)$) to a family of compact dualizable generators (the powers of $\omega_{\G/\Mfg}$), we may apply \Cref{prop:rigid} to obtain the following lemma:

\begin{lem}\label{lem:MU-koszul}
  There is an equivalence of presentably symmetric monoidal categories
  \[ \mathrm{IndCoh}(\mathcal{M}_{\mathrm{fg}}) \simeq \mathrm{Mod}(\Ab^{\Gr}; p_*\mathcal{O}_{\Mfg} ). \]
\end{lem}

%% \begin{proof}
%%   By \Cref{lem:omega-euler-one}, there is a symmetric monoidal functor $\Z \to \IndCoh(\Mfg) ^{\heartsuit}$ sending $1 \in \Z$ to $\omega_{\G/\Mfg}$.
  
%%   As a consequence, we obtain a symmetric monoidal left adjoint
%%   \[\Ab^{\Gr} \to \mathrm{IndCoh}(\mathcal{M}_{\mathrm{fg}}) \]
%%   which sends $\Z(n)$ to $\omega^{n} _{\mathbb{G}/\Mfg}$.
%%   This functor sends a family of compact dualizable generators to a family of compact dualizable generators,
%%   so by [ref] \todo{Find or write ref!} it factors through the desired symmetric monoidal equivalence
%%   \[  \mathrm{Mod}(\Ab^{\Gr}; \ExtMUt ) \simeq \mathrm{IndCoh}(\mathcal{M}_{\mathrm{fg}}). \qedhere \]
%%   %% \[ \mathrm{Mod}(\Ab^{\Gr}; \ExtMU) \to \mathrm{IndCoh}(\mathcal{M}_{\mathrm{fg}}). \]
%% \end{proof}

Using \Cref{lem:MU-koszul} and \Cref{lem:tensor-mod}, we obtain the following corollary:

\begin{cor}\label{lem:uMU-koszul}
  There is an equivalence of presentably symmetric monoidal categories
  \[ \uAb_{i2} \otimes_{\Ab} \mathrm{IndCoh}(\mathcal{M}_{\mathrm{fg}}) \simeq \Mod ( \uAb_{i2}^{\Gr}; \underline{p_*\mathcal{O}_{\Mfg}} \otimes_{\uZ} \uZt ).  \]
\end{cor}

%% \begin{proof}
%%   Using \Cref{lem:MU-koszul}, we have equivalences
%%   \begin{align*}
%%     \uAb \otimes_{\Ab} \mathrm{IndCoh}(\mathcal{M}_{\mathrm{fg}})
%%     &\simeq \uAb \otimes_{\Ab} \mathrm{Mod}(\Ab^{\Gr}; \ExtMUt ) \\
%%     &\simeq \uAb^{\Gr} \otimes_{\Ab^{\Gr}} \mathrm{Mod}(\Ab^{\Gr}; \ExtMUt ).
%%   \end{align*}
%%   We conclude by applying \Cref{lem:tensor-mod}.
%% \end{proof}



\subsection{The main lemma} \label{subsec:main-lem}\ 


In this subsection, we prove our main lemma, \Cref{lem:main-lem}, and use it to prove Theorems \ref{thm:Cta} and \ref{thm:Cta-homotopy}. This lemma gives us an explicit formula for $\Gr(i_*(\Gamma_*(-)))$ on a restricted class of objects. Once we have this formula the remaining work is relatively easy.
We begin with a definition and a couple of useful lemmas.

\begin{dfn}\ 
  \begin{itemize}
  %\item Let $\Sp_{C_2, i2} ^{\MURev}$ denote the full subcategory of $E \in \Sp_{C_2, i2}$
  %  for which $\MURu_{n\rho - 1} ^{C_2} E = 0$ for all $n$.
  \item Let $\Sp_{C_2, i2} ^{\mathrm{proj}}$ denote the full subcategory of $E \in \Sp_{C_2, i2}$ for which $\MU_{\R,2} \otimes E$ is a retract of a sum of pure suspensions of $\MU_{\R,2}$ (a suspension $\Sigma^{a+b\sigma} \MU_{\R,2}$ is said to be pure if $a=b$). 
    %% $C_2$-spectrum of the form $\bigoplus_{\alpha} \Sigma^{n_\alpha \rho} \MUR$ for some collecion $n_\alpha$ of integers.
  \item Let $\Sp_{i2} ^{\mathrm{proj}}$ denote the full subcategory of $E \in \Sp_{i2}$ for which $\MU_{2} \otimes E$ is a retract of a sum of even suspensions of $\MU_2$.
  \end{itemize}
  Since the underlying spectrum of $\MU_{\R, 2}$ is $\MU_2$, the underlying spectrum of an object of $\Sp_{C_2, i2} ^{\mathrm{proj}}$ is contained in $\Sp_{i2} ^{\mathrm{proj}}$.
  Note that both $\Sp_{C_2,i2} ^{\mathrm{proj}}$ and $\Sp_{i2} ^{\mathrm{proj}}$ contain units and are closed under tensor product, so each inherits the structure of a symmetric monoidal category.
\end{dfn}

%% \begin{rec}
%%   There is an equivalence of Hopf algebroids in $\uAbh$ between $(\underline{\MUR}_{\rho *}, \underline{\MUR}_{\rho *} \MUR)$ and $\underline{(\MU_*, \MU_*\MU)}$.
%% \end{rec}


%% \begin{cnstr}
%%   We construct a homology functor $h$
%% \end{cnstr}

%% \begin{lem}
%%   The homology functor $h$ can be equipped with a symmetric monoidal structure which makes the following diagram of lax monoidal functor commute,
%%   \begin{center} \begin{tikzcd}
%%       \Sp_{C_2} \ar[d, "h"] \ar[r, " - \otimes \mathrm{cb}(\MU_\R)"] &
%%       \Sp_{C_2}^{\Delta} \ar[r, "\piu_0^{\rho}Y"] &
%%       \Sp_{C_2}^{\Gr, \rho, \heartsuit, \Delta} \ar[r, "\mathrm{Tot}"] &
%%       \Sp_{C_2}^{\Gr} \ar[d] \\
%%       (\uAb \otimes_\Z \IndCoh (\Mfg))^{\heartsuit} \ar[r, "p_*"] &
%%       \uAb^{\Gr} \ar[r, "\twist^{\rho}"] &
%%       \uAb^{\Gr} \ar[r] &
%%       \Sp_{C_2}^{\Gr}
%%   \end{tikzcd} \end{center}
%% \end{lem}

%% \begin{proof}

%% \end{proof}



%\begin{lem} \label{lem:MUR-uZ-mod}
  %Let $E \in \Sp_{C_2} ^{\MURev}$. Then the restriction map on $\MURu_{n\rho} ^{C_2} E$ is injective, which implies that $\MURu_{n\rho} ^{C_2} E$ is a $\uZ$-module.
%\end{lem}
%
%\begin{proof}
  %Looking at maps out of the cofiber sequence
  %\[(C_2)_+ \otimes \Ss^{n\rho} \to \Ss^{n\rho} \to \Ss^{n\rho+\sigma}\]
  %implies that the kernel of the restriction map $(\MUR)_{n\rho} E \to \MU_{2n} \Phi^e(E)$ is equal to the image of a map $(\MUR)_{n\rho+\sigma} E \to (\MUR)_{n\rho} E$.
  %The domain is zero by assumption.
  %To show that $\MURu_{n\rho} E$  is a $\uZ$-modules, we must show that $[C_2] = 2$.
  %This follows from the relation $ r \circ [C_2] = r \circ 2 $
  %% \[ r \circ [C_2] = r \circ t \circ r = (1+\sigma) \circ r = 2 \circ r = r \circ 2\]
  %and the fact that $r$ is injective as shown above.
%\end{proof}

\begin{lem} \label{lem:MUR-uZ-mod}
  Suppose that $E \in \Sp_{C_2, i2} ^{\mathrm{proj}}$. Then 
  \[\underline{(\MU_{\R,2})}_{n\rho} ^{C_2} (E) \cong \underline{(\MU_2)_{2*} (E)}.\]
  As a consequence, $\underline{(\MU_{\R,2})}_{n\rho} ^{C_2} (E)$ acquires the structure of a $\uZt$-module.
\end{lem}

\begin{proof}
  By definition of $\Sp_{C_2, i2} ^{\mathrm{proj}}$, it suffices to note that this is true for $E = \Ss_2$, which follows from \cite[Theorem 2.28]{HuKriz}.
\end{proof}

\begin{lem} \label{lem:main-lem}
  Let $h$ denote the symmetric monoidal functor $h : \Sp_{i2} ^{\mathrm{proj}} \to \IndCoh(\Mfg)_{i2} ^{\heartsuit}$
  which sends a spectrum to its associated $((\MU_2)_*, (\MU_2)_*\MU_2 )$-comodule.  
  There is a commutative diagram of lax symmetric monoidal functors
  \begin{center}
    \begin{tikzcd}
      \Sp_{C_2, i2} ^{\mathrm{proj}} \ar[rrrr,"i_* \circ \Gamma_*"] \ar[d, "\Phi^e"] & & & &
      \Sp_{C_2, i2} ^{\Fil} \ar[d, "\Gr"] \\
      \Sp^{\mathrm{proj}} _{i2} \ar[r, "h"] &
      \IndCoh(\Mfg)_{i2} ^{\heartsuit} \ar[r, "p_*"] &
      \Ab^{\Gr}_{i2} \ar[r, "\_"] &
      \uAb^{\Gr}_{i2} \ar[r, "\twist^\rho"] &
      \Sp_{C_2,i2} ^{\Gr}.
    \end{tikzcd}
  \end{center}
  %% where we have postcomposed $\twist^{\rho}$ with the forgetful functor from $\uAb^{\Gr}$ to $\Sp_{C_2} ^{\Gr}$, and $F$ is the functor $F(E) = \Tot^\bullet (\MURu_{*\rho} E \otimes_{\MURu_{*\rho}} \MURu_{*\rho} \MUR ^{\otimes \bullet +1})$ taking $E$ to the cobar complex of $\MURu_{*\rho} E$ considered as a $\MURu_{*\rho} \MUR$-comodule.

  %% Moreover, $F$ lands in the full subcategory of $\uAb^{\Gr}$ for which the the restriction $r$ is an equivalence, and there is a further diagram
  %% \begin{center}
  %%   \begin{tikzcd}
  %%     \Sp_{C_2} ^{\MURpr} \ar[rr,"\Gamma_*"] \ar[d, "G"] &  & \Sp_{C_2} ^{\Fil} \ar[d, "\Gr"] \\
  %%     \Ab^{\Gr} \ar[r, "\_"] & \uAb^{\Gr} \ar[r, "\twist^\rho"] & \Sp_{C_2} ^{\Gr},
  %%   \end{tikzcd}
  %% \end{center}
  %% where $G$ is the functor $G(E) = \Tot^{\bullet} (\MU_{2*} E \otimes _{\MU_{2*}} \MU_{2*} \MU^{\otimes \bullet+1})$ taking $E$ to the cobar complex of $\MU_{2*} E$ considered as a $\MU_{2*} \MU$-comodule.
\end{lem}

\begin{proof}
  %% Let $\Mod(\Sp_{C_2} ^{\Delta}; \MUR^{\bullet + 1})^{\mathrm{pt-proj}}$ denote the full symmetric monoidal subcategory of $\Mod(\Sp^{\Delta}; \MUR^{\bullet + 1})$ consisting of the objects which are pointwise direct summands of $\bigoplus_\alpha \Sigma^{n_{\alpha} \rho} \MUR^{k+1}$ as $\MUR^{k+1}$-modules.

  %% This is defined so that the $\MUR$-cobar complex defines a symmetric monoidal functor
  %% \[\Sp_{C_2} ^{\MURpr} \to \Mod(\Sp_{C_2} ^{\Delta}; \MUR^{\bullet + 1})^{\mathrm{pt-proj}}.\]

  %% We claim that the following diagram of lax symmetric monoidal functors commutes:

  Our argument will rest on the existence of the following commuting diagram of lax symmetric monoidal functors.
  \begin{center}
    \begin{tikzcd}
      & \Sp_{i2} ^{\mathrm{proj}} \ar[ddl, "h"] \ar[d, "- \otimes \mathrm{cb}(\MU_2)"] & &
      \Sp_{C_2,i2} ^{\mathrm{proj}} \ar[d, "- \otimes \mathrm{cb}(\MU_{\R,i2})"] \ar[ddl, dashed, "f"'] \ar[ll, "\Phi^e"] \\
      & \Sp_{i2}^{\Delta} \ar[d, "\piu_0^2Y"] & &
      \Sp_{C_2,i2} ^{\Delta} \ar[ll, "\Phi^e"] \ar[r,"\tau_{\geq 0} ^\rho Y"'] \ar[d, "\piu_{0}^{\rho} Y"'] &
      \Sp_{C_2,i2} ^{\Fil, \Delta} \ar[d, "\Gr"] \\
      \IndCoh(\Mfg)_{i2}^{\heartsuit} \ar[ddr, "p_*"] & \Ab_{i2}^{\Gr,2,\heartsuit,\Delta} \ar[d, "\twist^{-2}"] &
      \uAb_{i2} ^{\Gr, \rho,\heartsuit, \Delta} \ar[l, "\Phi^e"] \ar[d, "\twist^{-\rho}"] \ar[r, "\mathrm{forget}"] &
      \Sp_{C_2,i2} ^{\Gr, \rho, \heartsuit, \Delta} \ar[r] & \Sp_{C_2,i2}^{\Gr, \Delta} \ar[dd, "\Tot"] \\
      & \Ab_{i2}^{\Gr,\heartsuit,\Delta} \ar[d, "\Tot"] &
      \uAb_{i2}^{\Gr, \heartsuit, \Delta} \ar[l, "\Phi^e"] \ar[d, "\Tot"] & & \\
      & \Ab_{i2}^{\Gr} &
      \uAb_{i2}^{\Gr} \ar[r, "\twist^{\rho}"] \ar[l, "\Phi^e"] &
      \uAb_{i2}^{\Gr} \ar[r,"\mathrm{forget}"] &
      \Sp_{C_2, i2} ^{\Gr}.
    \end{tikzcd}
  \end{center}

  The squares formed by the underlying functor, $\Phi^e$, commute since underlying commutes with limits and colimits, the underlying of $\Ss_2^\rho$ is $\Ss_2^2$ and the underlying of $\MU_{\R,2}$ is $\MU_2$.
  In grid position $(4,3)$ we have implicitly made use of the equivalence $\Sp_{C_2} ^{\Fil,\rho,\heartsuit,\Delta} \simeq \Sp_{C_2} ^{\Gr, \rho, \heartsuit, \Delta}$ to identify the target of $\piu_0 ^{\rho}Y$ with the latter. The upper-right square commutes by \Cref{prop:twist-slice}. The dashed arrow, labelled $f$, is unique (and lax symmetric monoidal) if it exists since the forgetful functor is fully-faithful. The dashed arrow then exists by \Cref{lem:MUR-uZ-mod}. The large bottom-right square commutes because $\twist^{\rho}$ commutes with $\Tot$. Finally, the existence of the factorization along the left side is a corollary of our ability to identify the $\mathrm{E}_2$-page of the Adams--Novikov spectral sequence (see \cite{Adams}). %(see \cite{Pstragowski}).
  %\todo{Seriously? See Piotr for this?}

  Using the fact that $\Gr$ and $\Tot$ commute, we have equivalences
  \[ \Gr \circ i_* \circ \Gamma_* \simeq \Tot  \Gr ( \tau_{\geq 0}^\rho Y ( - \otimes \mathrm{cb}(\MU_{\R,2}))) \simeq \twist^{\rho} (\Tot (\twist^{-\rho}( f(-)))). \]
  By \Cref{lem:MUR-uZ-mod}, the functor $f$ (and its composition with the twist) lands in the full subcategory of the target which is spanned, in each grading and cosimplicial degree, by $\uZt$-modules for which the restriction map $r$ is an equivalence. Thus, we have an equivalence of lax symmetric monoidal functors
  \[ \twist^{\rho} (\Tot (\twist^{-\rho}( f(-)))) \simeq \twist^{\rho} (\Tot (\underline{\Phi^e \twist^{-\rho}( f(-))})). \]
  Since $\_$ commutes with limits we then have further equivalences,
  \[ \twist^{\rho} (\Tot (\underline{\Phi^e \twist^{-\rho}( f(-))}))
  \simeq \twist^{\rho} ( \underline{\Tot (\Phi^e \twist^{-\rho}( f(-))) } )
  \simeq \twist^{\rho} ( \underline{  p_* h \Phi^e(-) } ). \qedhere \]
  %% \begin{center}
  %%   \begin{tikzcd}
  %%     \Sp_{C_2} ^{\MURpr} \ar[rr, "\MUR\mathrm{-cobar}"] \ar[rddd, bend right, "F"'] & & \Mod(\Sp_{C_2} ^{\Delta}; \MUR^{\bullet + 1})^{\mathrm{pt-proj}} \ar[r,"\tau_{\geq 0} ^\rho"] \ar[d, "\piu_{0}^{\rho}"] \ar[ld, "\{\Sigma^{n\rho} \piu_{n\rho} ^{C_2}\}"'] & \Sp_{C_2} ^{\Fil, \Delta} \ar[d, "\Gr"] \\
  %%      & \uAb^{\Gr, \rho,\heartsuit, \Delta} \ar[d, "\twist^{-\rho}"] \ar[r, "\mathrm{forget}"] & \Sp_{C_2} ^{\Gr, \rho, \heartsuit, \Delta} \ar[r] & \Sp_{C_2}^{\Gr, \Delta} \ar[dd, "\Tot"] \\
  %%      & \uAb^{\Gr, \heartsuit, \Delta} \ar[d, "\Tot"] & & \\
  %%      & \uAb^{\Gr} \ar[r, "\twist^{\rho}"] & \uAb^{\Gr} \ar[r,"\mathrm{forget}"] & \Sp_{C_2} ^{\Gr}
  %%   \end{tikzcd}
  %% \end{center}
  
  %%   By \Cref{prop:twist-slice} and the fact that $\Gr$ and $\Tot$ commute, we see that the desired commutative square may be obtained by going around the outside of our diagram. It therefore suffices to prove that the above diagram commutes.

  %% The leftmost region of the diagram commutes by the definition of $F$.
  

  %% , it follows that the composite \[\{\piu_{n\rho} ^{C_2}\} : \Mod(\Sp_{C_2} ^{\Delta}; \MUR^{\bullet + 1})^{\mathrm{pt-proj}} \xrightarrow{\{\Sigma^{n\rho} \piu_{n\rho} ^{C_2}\}} \uAb^{\Gr, \rho, \heartsuit, \Delta} \xrightarrow{\twist^{-\rho}} \uAb^{\Gr, \heartsuit, \Delta}\]
  

  %% By \Cref{prop:const-mackey}, the composite $\uAb^{\Gr, \heartsuit, \Delta} \xrightarrow{\mathrm{forget}} \Ab^{\Gr, \heartsuit, \Delta} \xrightarrow{\_} \uAb^{\Gr, \heartsuit, \Delta}$ is naturally ismorphic to the identity.

  %% Using the fact that $\_$ and $\mathrm{forget}$ commute with totalizations,\todo{Probably should have recorded this earlier for $\_$... what's the simplest way to get this?} we obtain the second commuting square.
\end{proof}

Finally, we are able to prove our main theorems:

\begin{proof}[Proof of \Cref{thm:Cta} and \Cref{thm:Cta-homotopy} (part 1).]
  We have a chain of symmetric monoidal equivalences:

  \begin{align*}
    \mathrm{Mod}(\SHRat; C\ta)
    &\xrightarrow[\ref{lem:gr-of-comparison}]{\simeq} \mathrm{Mod}( \Sp_{C_2, i2}^{\Gr}; \Gr(R_{\bullet}) ) 
    \xrightarrow[\ref{lem:main-lem}]{\simeq} \mathrm{Mod}( \uAb_{i2}^{\Gr}; \twist^{\rho} (\underline{p_*\mathcal{O}_{\Mfg}}) ) \\
    &\xrightarrow[\twist^{-\rho}]{\simeq} \mathrm{Mod}( \uAb_{i2}^{\Gr}; \underline{p_*\mathcal{O}_{\Mfg}} ) 
    \xrightarrow[\ref{lem:uMU-koszul}]{\simeq} \uAb_{i2} \otimes_{\Ab} \IndCoh(\Mfg)
  \end{align*}

  The claim about the $\Sp_{C_2,i2}$-algebra structure follows once we know $\Sp_{C_2,i2}$ acts through the ambient category at each step. 
  The key points here are
  that the equivalence from \Cref{lem:main-lem} was induced by an equivalence of algebras and
  that the twist functor is an $\uAb_{i2}$-algebra map.

  Using the $\Sp_{C_2,i2}$-linearity, the claim about Picard elements reduces to tracking $C\ta \otimes \Ss_2^{0,0,1}$ around the diagram. Under the first equivalence, $C\ta \otimes \Ss_2^{0,0,1}$ goes to $\o (1)$. The twist by $-\rho$ sends this to $\Sigma ^{\rho}\o (1)$. The final equivalence sends $\o(1)$ to $\uZt \otimes \omega_{\mathbb{G}/\Mfg}$. Altogether we learn that $C\ta \otimes \Ss_2^{0,0,1}$ is sent to $\Sigma^{-\rho}\uZ_2 \otimes \omega_{\mathbb{G}/\Mfg}$, as desired.

  Now suppose that $X \in \Sp_{C_2}^{\mathrm{proj}}$.
  Then, using \Cref{lem:main-lem}, we have
  \[\Gr(i_* (\Gamma_* X)) \simeq \twist^{\rho} (\underline{p_* h \Phi^e (X) }). \]
  Continuing along the sequence of equivalences above, we see that this object is sent to
  $\uZt \otimes h(\Phi^e(X))$, as desired. \qedhere

  %% We learn that, under the equivalence $\Mod(\SHR^{\at}; C\ta) \simeq \Mod(\uAb ^{\Gr}; \underline{\ExtMUt})$, $\Gamma_* (X) \otimes C\ta$ corresponds to $\underline{C^* _{\MU_{2*} \MU} (\MU_{2*} (X))}$.
  %% Under the equivalence $\Mod(\uAb ^{\Gr}; \underline{\ExtMUt}) \simeq \uAb \otimes_\Ab \IndCoh(\Mfg)$, this corresponds to the comodule $\underline{\MU_{2*} X} \cong \MURu_{*\rho} X$.

  
  %% From \Cref{lem:gr-of-comparison}, we have an equivalence of symmetric monoidal categoies,
  %% \[ \mathrm{Mod}(\SHRat_{i2}; C\ta) \simeq \mathrm{Mod}( \Sp_{C_2, i2}^{\Gr}; \Gr(\Gamma_*\Ss_2) ). \] 
  %% Using \Cref{lem:main-lem}, we have an equivalence of commutative algebra,
  %% \[\Gr(\Gamma_* \Ss_2) \xrightarrow[\ref{lem:gr-of-comparison}]{\simeq} \twist^{\rho} (\underline{p_*\O_{\Mfg}}).\]
  
  %% %% By definition, we have $G(\Ss_2) \simeq \ExtMUt$. It follows that we have equivalences of symmetric monoidal categories,
  %% %% \[ \Mod (\SHR^{\at}; C\ta) \simeq \Mod(\uAb ^{\Gr}; \twist^{\rho} (\underline{\ExtMUt})). \]
  %% Finally, applying $\twist^{-\rho}$, we obtain a symmetric monoidal equivalence
  %% \[\Mod(\uAb ^{\Gr}; \twist^{\rho} (\underline{\ExtMUt})) \simeq \Mod(\uAb ^{\Gr}; \underline{\ExtMUt}).\]
  %% By \Cref{lem:uMU-koszul}, this is equivalent to $\uAb \otimes_\Ab \IndCoh(\Mfg)$, as desired.

  %% %% \begin{align*}
  %% %%   \Mod (\SHR^{\at}; C\ta) &\simeq \Mod(\Sp_{C_2} ^{\Gr}; \Gr(\Gamma_* \Ss_2)) \\
  %% %%   &\simeq \Mod(\Sp_{C_2} ^{\Gr}; \twist^{\rho} (\underline{\ExtMUt})) \\
  %% %%   &\simeq \Mod(\uAb ^{\Gr}; \twist^{\rho} (\underline{\ExtMUt})).
  %% %% \end{align*}
  %% %% Finally, applying $\twist^{-\rho}$, we obtain a symmetric monoidal equivalence
  %% %% \[\Mod(\uAb ^{\Gr}; \twist^{\rho} (\underline{\ExtMUt})) \simeq \Mod(\uAb ^{\Gr}; \underline{\ExtMUt}).\]
  %% %% By \Cref{lem:uMU-koszul}, this is equivalent to $\uAb \otimes_\Ab \IndCoh(\Mfg)$, as desired.
\end{proof}

In order to calculate the homotopy groups of $C\ta$-modules, we begin with a simple lemma.

\begin{lem} \label{lem:underline-homotopy}
  Given any $A \in \Ab$, we have
  \[\pi_{p+q\sigma} ^{C_2} \underline{A} \cong \bigoplus_{a} \pi_{p-a} \left(A \otimes \pi_{a+q\sigma} \uZ \right).\]
\end{lem}

\begin{proof}
  Consider the composite
  $ \Ab \to \uAb \to \Ab^{\Gr} $
  given by
  $ A \mapsto \underline{A} \mapsto \{\Map(\Ss^{q\sigma}, \underline{A})\} $.
  This composite is colimit-preserving, hence is equivalent to
  $ A \mapsto \{A \otimes \Map(\Ss^{q\sigma} , \underline{\Z})\} $.
  Now, we have
  \[\Map(\Ss^{q\sigma} , \underline{\Z}) \simeq \bigoplus_a \Sigma^{a} \pi_{a+q\sigma} \uZ \]
  because it is a $\Z$-module, and the result follows by applying $\pi_p$.
\end{proof}

\begin{proof}[Proof of \Cref{thm:Cta} and \Cref{thm:Cta-homotopy} (part 2).]
  We now turn to our assertions about homotopy groups.
  Tracing through the various equivalences, we find that for $X$ which is $\MU_{\R,2}$-projective
  \[ \pi_{p,q,w}^\R(C\ta \otimes X) \cong \pi_{(p-w) + (q-w)\sigma}^{C_2} ( (\uZt \otimes_{\Zt} p_*h \Phi^e (X))_w ). \]
  We can commute taking the $w^{\mathrm{th}}$ component past the tensor product and then apply \Cref{lem:underline-homotopy} and the fact that $p_*h$ computes $\Ext_{(\MU_2)_*\MU_2}$ to conclude that
  \begin{align*}
    \pi_{p,q,w}^\R(C\ta \otimes X)
    &\cong \pi_{(p-w) + (q-w)\sigma}^{C_2} (\uZt \otimes_{\Zt} (p_*h \Phi^e (X))_w ) \\
    &\cong \bigoplus_a \pi_{p-w-a} (\pi_{a+(q-w)\sigma}\uZt \otimes_{\Zt} \Ext_{(\MU_2)_*\MU_2}^{*,2w}((\MU_2)_*, (\MU_2)_* \Phi^e (X) )) \\
    &\cong \bigoplus_a \pi_{p-w-a} (\Ext_{(\MU_2)_*\MU_2}^{*,2w}((\MU_2)_*, \pi_{a+(q-w)\sigma}\uZt \otimes_{\Zt} (\MU_2)_* \Phi^e (X) )) \\
    &\cong \bigoplus_a \Ext_{(\MU_2)_*\MU_2}^{a+w-p,2w}((\MU_2)_*, \pi_{a+(q-w)\sigma}\uZt \otimes_{\Zt} (\MU_2)_* \Phi^e (X) ). 
  \end{align*}
  Since $\pi_{a+q\sigma} ^{C_2} \uZt$ is isomorphic to $\Z_2$ or $\F_2$, and $(\MU_2)_* X$ is torsion-free (since it is a projective $(\MU_2)_*$-module), the tensor product inside the Ext can be taken in a $1$-categorical sense.
  %
  %% the place where the homotopy groups are most accessible is in the category $\uAb^{\Gr}$ where we are thinking about the object $\uZt \otimes p_*h \Phi^e (X)$. 
  %
  %% Finally, we prove the 
  %% Suppose that $E \in \Mod(\SHR^{\at}; C\ta)$ corresponds to $\{M_w\} \in \Mod(\uAb^{\Gr}; \underline{\ExtMUt})$ under the equivalence defined above.
  %% Then, tracing through the above equivalences of categories, we find that
  %% \[\pi_{p,q,w} ^\R E \cong \pi_{(p-w) + (q-w)\sigma}^{C_2} M_w.\]
  %% When $E = \Gamma_* (X) \otimes C\ta$ for $X \in \Sp_{C_2} ^{\MURpr}$, we find that
  %% \[M_w \simeq \underline{\Ext^{*, 2w} _{\MU_* \MU} (\MU_* X)},\]
  %% where
  %% \[\pi_p \Ext^{*, 2w} _{\MU_* \MU} (\MU_* X) \cong \Ext^{-p, 2w} _{\MU_* \MU} (\MU_* X).\]
  %% By \Cref{lem:underline-homotopy}, we have
  %% \[\pi_{(p-w)+(q-w) \sigma} ^{C_2} \underline{\Ext^{*, 2w} _{\MU_* \MU} (\MU_* X)} \cong \bigoplus_{a} \pi_{p-w-a} (\Ext^{*, 2w} _{\MU_* \MU} (\MU_* X) \otimes \pi_{a+(q-w)\sigma} ^{C_2} \uZ).\]
  %% Since $\pi_{a+q\sigma} ^{C_2} \Z$ is isomorphic to $\Z$ or $\F_2$, and $\MU_* X$ is torsion-free since it is a projective $\MU_*$-module, we find that 
  %% \[\Ext^{*, 2w} _{\MU_* \MU} (\MU_* X) \otimes \pi_{a+(q-w)\sigma} ^{C_2} \uZ \simeq \Ext^{*, 2w} _{\MU_* \MU} (\MU_* X \otimes \pi_{a+(q-w)\sigma} ^{C_2} \uZ), \]
  %% from which we deduce that
  %% \begin{align*}
  %%   \pi_{p,q,w} ^{\R} \Gamma_* (X) \otimes C\ta &\cong \bigoplus_{a} \pi_{p-w-a} \Ext^{*, 2w} _{\MU_* \MU} (\MU_* X \otimes \pi_{a+(q-w)\sigma} ^{C_2} \uZ) \\
  %%   &\cong \bigoplus_{a} \Ext^{a+w-p, 2w} _{\MU_* \MU} (\MU_* X \otimes \pi_{a+(q-w)\sigma} ^{C_2} \uZ) \\
  %%   &\cong \bigoplus_{w+a-s=p} \Ext^{s, 2w} _{\MU_* \MU} (\MU_* X \otimes \pi_{a+(q-w)\sigma} ^{C_2} \uZ),
  %% \end{align*}
  %% as desired.
\end{proof}
%
%\begin{lem}\label{lem:gr-R-bullet-computation}
%  The graded $\mathbb{E}_\infty$-ring in $\Sp_{C_2,i2}$ given by $\Gr(\Gamma_* (\Ss_2))$ is equivalent to the image of $\ExtMUt$ under the composite $\twist^{\rho} \circ \_ : \Ab^{\Gr} \to \uAb^{\Gr} \to \CAlg(\Sp_{C_2}^{\Gr})$.
%\end{lem}\todo{this needs a twist}
%
%The main idea of the proof is simple, and not well represented due to technical details.
%Upon examining the totalization computing $\Gr(R_\bullet)$ given in [location] directly, one finds that the full subcategory of $\Sp_{C_2}^{\Gr}$ on the relevant objects is a 1-category. This ``collapse of coherences'' allows us to factor the cosimplicial diagram through the right adjoint $- \otimes_{\Z} \uZt$ by hand.
%
%\begin{proof}  
  %% This lemma is the heart of the content of this section.
  %% For this reason, here and only here, we will need to examine the construction of $R_\bullet$ directly.
%
%  We begin by recalling the construction of $R_\bullet$,
%  \[ R_\bullet \coloneqq \mathrm{Tot}^* \left( P_{2 \bullet} \MU_\R^{\otimes * + 1} \right). \]
%  Using the material on twisted $t$-structures we can rewrite this as follows,
%  \[ R_\bullet \coloneqq \mathrm{Tot}^* \tau_{\geq 0}^{\rho} \mathrm{const}_\bullet (\MU_\R^{\otimes * + 1} ). \]
%  The operation of passing to associated graded commutes with limits and colimits so we have an equivalence of graded $\mathbb{E}_\infty$-rings in $C_2$-spectra,
%  \[ \Gr_\bullet \mathrm{Tot}^* \tau_{\geq 0}^{\rho} \mathrm{const}_\bullet (\MU_\R^{\otimes * + 1} ) \simeq \mathrm{Tot}^* \Gr_\bullet \tau_{\geq 0}^{\rho} \mathrm{const}_\bullet (\MU_\R^{\otimes * + 1} ). \]
%  Since $\piu_0^{\rho}$ kills the shift map, there is a natural ring map,\todo{Needs work}
%  \[ \Gr_\bullet \tau_{\geq 0}^{\rho} \mathrm{const}_\bullet (\MU_\R^{\otimes * + 1} ) \to \piu_{0}^{\rho} \mathrm{const}_\bullet (\MU_\R^{\otimes * + 1} ) \]
%  which from our knowledge of the slices of $\MU_\R$ we may conclude is an equivalence.
%  This provides a factorization,
%  \[ \Delta \xrightarrow{F_{\bullet, *}} \mathrm{CAlg}( \uAb^{\Gr, \heartsuit, \rho} ) \simeq \mathrm{CAlg}(\Sp_{C_2}^{\Gr, \heartsuit, \rho}) \to \mathrm{CAlg}(\Sp_{C_2}^{\Gr}) \]
%  of the diagram whose totalization we would like to understand.
%  We can also express $\ExtMU$ as the totalization of a cosimplicial diagram which factors through commutative algebras in the heart---the $\MU$ cobar complex. Thus, in order to prove the lemma it will suffice to show that the left triangle in the diagram below commutes.
%
  %% \begin{center}
  %%   \begin{tikzcd}
  %%     & \mathrm{CAlg} ( \Ab^{\Gr, \heartsuit}) \ar[d] \ar[r] & \mathrm{CAlg} ( \Ab^{\Gr} ) \ar[d] \\
  %%     & \mathrm{CAlg} ( \uAb^{\Gr, \heartsuit}) \ar[d, "twist"] \ar[r] & \mathrm{CAlg} ( \uAb^{\Gr} ) \ar[d, "twist"] \\
  %%     \Delta \ar[uur, "MU cobar"]  \ar[r, "F_{*,\bullet}"]
  %%     & \mathrm{CAlg} ( \Ab^{\Gr, \heartsuit, \rho}) \ar[d] \ar[r] & \mathrm{CAlg} ( \Ab^{\Gr} ) \ar[d] \\      
  %%   \end{tikzcd}
  %% \end{center}
%
%  \begin{center}
%    \begin{tikzcd}
%      & & \Delta \ar[ddll, "\MU\text{cobar}"] \ar[dd, "F_{*,\bullet}"] \ar[ddr, "\MU \text{cobar}"] \\ \\
%      \mathrm{CAlg} ( \Ab^{\Gr, \heartsuit} ) \ar[r, "- \otimes_{\Z} \uZ"] &
%      \mathrm{CAlg} ( \uAb^{\Gr, \heartsuit} ) \ar[r, "\text{twist}"] &
%      \mathrm{CAlg} ( \uAb^{\Gr, \heartsuit, \rho} ) \ar[r, ] &
%      \mathrm{CAlg} ( \Ab^{\Gr, \heartsuit} )
%    \end{tikzcd}
%  \end{center}
%
%  We begin by observing that the big triangle with two legs given by the $\MU$ cobar complex commutes since the horizontal composite accross is the identity.
%  Since $U \piu_0^\rho(-) \cong \piu_0^{S^2}(U-)$, the right triangle commutes by the definition of the $\MU$ cobar complex.
%  Now we make an observation that is special to $F_{*,\bullet}$. 
%  Since the slices of $\MU_\R^{\otimes * + 1}$ are all of the form $\Sigma^{n \rho}\uZt$ the underlying non-equivariant object functor is fully faitful when restricted to relevant objects [ref]. This implies that the 2-out-of-3 property applies to the commutativity of this diagram, so we may conclude the left triangle commutes, as desired.
%
%% AAAAAAAAAAAAAAAAAAAAAAAAAAAAAAAAA
%
%  
%%   Recall that $R_\bullet$ is constructed as the totalization of the cosimplicial CAlg in filtered $C_2$-spectra below,\footnote{Here and below we will use decorations $\bullet$ and $*$ to disambiguate indices.}
%%   \[ R_\bullet \coloneqq \mathrm{Tot}^* \left( P_{2 \bullet} \MU_\R^{\otimes * + 1} \right). \]  
%%   The operation of passing to associated graded commutes with limits and colimits so we have an equivalence of graded $\mathbb{E}_\infty$-rings in $C_2$-spectra,
%%   \[ \Gr_\bullet \mathrm{Tot}^* \left( P_{2\bullet} \MU_\R^{\otimes * + 1} \right) \simeq \mathrm{Tot}^* \left( \Gr_\bullet P_{2\bullet} \MU_\R^{\otimes * + 1} \right). \]
%%   Using the eveness of $\MU_\R$ we know that for fixed values of $\bullet$ and $*$ these are all just sums of copies of $\Sigma^{n\rho} \uZt$. Moreover, this means all the maps are determined by what they do on the homotopy groups of the underlying spectrum, where we can drop all the $\R$'s. At this level we're just looking at the $\MU$-cobar complex. However, since we need to understand the commutative algebra structure in a very precise way we take more time in coming to this point.
%
%  
%%   In [location] we constructed $R_\bullet$ with its ring structure as follows,
%%   \[  R_\bullet \coloneqq \mathrm{Tot}^* \tau_{\geq 0}^{\rho} \mathrm{const} ( \MU_\R^{\otimes * + 1} ), \]
%%   where const refers to the constant filtered object functor.
%  
%  
%
%%   \[  \Gr_\bullet P_{2\bullet} \MU_\R^{\otimes * + 1} \simeq  P_{2\bullet}^{2\bullet+1} \MU_\R^{\otimes * + 1}  \simeq  P_{2\bullet}^{2\bullet} \MU_\R^{\otimes * + 1}. \]
%%   In [location] we actually had a more tightly coupled version of this construction there $P_{2\bullet}$ was replaced with the composite of the constant tower functor with $\tau_{\geq 0}^{\mathrm{B}}$, the connective cover the for the slice version of the Beilinson $t$-structure on filtered objects. Using the fact that powers and slices of $\MU$ are so well behaved, we can rewrite things again as
%%   \[ \Gr R_\bullet \simeq \mathrm{Tot}^* \pi_0^{\mathrm{B}} \mathrm{const} ( \MU_\R^{\otimes * + 1} ). \]
%
%  
%%   \todo{prose here}
%  
%
%
%%   {\bf De-categorification}
%%   At this point we make an observation that greatly simplifies our understanding of $\Gr R_\bullet$:  
%%   Since the $\rho$-twisted $t$-structure on graded $C_2$-spectra is compatible with the tensor product, there is a pullback diagram of categories,
%%   \begin{center} \begin{tikzcd}
%%       \mathrm{CAlg}(\Sp_{C_2}^{\Gr,\heartsuit}) \ar[r] \ar[d] & \mathrm{CAlg}(\Sp_{C_2}^{\Gr}) \ar[d] \\
%%       \Sp_{C_2}^{\Gr,\heartsuit} \ar[r] & \Sp_{C_2}^{\Gr}
%%   \end{tikzcd} \end{center}
%%   where both horizontal arrows are fully faithful.
%%   Since, $F_{\bullet,*}$ lifts along the bottom horizontal arrow as a diagram of objects,
%%   we learn that it lifts along the upper arrow as a diagram of commutative algebras.
%%   The improvement here is that we're now working in a 1-category, therefore any sort of coherence statement becomes null. Thus, in order to produce the dashed lift indicated below, we only have to realize that we're working with objects which are sums of copies of $\uZt$.
%%   \begin{center} \begin{tikzcd}[sep = huge]
%%       &  \Delta \ar[ddr, "F_{\bullet,*}"] \ar[dl, dashed, bend right] & \\
%%       \mathrm{CAlg}(\Sp^{\Gr,\heartsuit}) \ar[r, hook] \ar[d, " - \otimes_{\Z} \uZ"] &
%%       \Z/\mathrm{CAlg}(\Sp^{\Gr}) \ar[d, " - \otimes_{\Z} \uZ"] & \\
%%       \mathrm{CAlg}(\Sp_{C_2}^{\Gr,\heartsuit}) \ar[r, hook] &
%%       (\underline{\pi}_0\o)/\mathrm{CAlg}(\Sp_{C_2}^{\Gr}) \ar[r] &
%%       \mathrm{CAlg}(\Sp_{C_2}^{\Gr})
%%   \end{tikzcd} \end{center}
  
%%   %% \begin{center} \begin{tikzcd}[sep = small]
%%   %%     & & & \Delta \ar[ddr, "F_{\bullet, *}"] \ar[dlll, dashed, bend right] & \\
%%   %%     \mathrm{CAlg}(\Sp^{\Gr,\heartsuit}) \ar[dd] \ar[rr, hook] \ar[dr, " - \otimes_{\Z} \uZ"] & &
%%   %%     \Z/\mathrm{CAlg}(\Sp^{\Gr}) \ar[dd] \ar[dr, " - \otimes_{\Z} \uZ"] \\
%%   %%     & \mathrm{CAlg}(\Sp_{C_2}^{\Gr,\heartsuit}) \ar[dd] \ar[rr, hook] & &
%%   %%     \underline{\pi}_0\o/\mathrm{CAlg}(\Sp_{C_2}^{\Gr}) \ar[r] \ar[dd] &
%%   %%     \mathrm{CAlg}(\Sp_{C_2}^{\Gr}) \ar[dd] \\
%%   %%     \Sp^{\Gr, \heartsuit} \ar[rr, hook] \ar[dr, " - \otimes_{\Z} \uZ"] & &
%%   %%     \mathrm{Mod}(\Sp^{\Gr}; \Z) \ar[dr, " - \otimes_{\Z} \uZ"] \\
%%   %%     & \Sp_{C_2}^{\Gr,\heartsuit} \ar[rr, hook] & &
%%   %%     \mathrm{Mod}(\Sp_{C_2}^{\Gr}; \underline{\pi}_0\o) \ar[r] &
%%   %%     \Sp_{C_2}^{\Gr} 
%%   %% \end{tikzcd} \end{center}
  
%%   In this diagram the right horizontal arrow preserves limits.
%%   We will argue that the right vertical arrow preserves limits as well so that we can evaluate the totalization in commutative algebras under $\Z$.
%%   Since passing to underlying is conservative and preserves limits it will suffice to show that
%%   \[ \mathrm{Mod}(\Sp^{\Gr}; \Z) \xrightarrow{- \otimes_{\Z} \uZ} \mathrm{Mod}(\Sp_{C_2}^{\Gr}; \underline{\pi}_0\o) \]
%%   preserves limits. It will suffice to argue this in the non-graded case.
%%   This functor can now be rewritten as the pullback functor
%%   \[ \mathcal{D}(\mathrm{Ab}) \to \mathcal{D}(\mathrm{Mackey}) \]
%%   induced by the map from the Burnside category to $\Z$ which gives $\uZ$. Since limits and colimits are computed pointwise we learn that this functor preserves all limits and colimits.

%%   Let $G_{\bullet,*}$ denote the cosimplicial diagram in commutative $\Z$-algebras which we wish to totalize. On the one hand the composite of $ - \otimes_{\Z} \uZ $ with the underlying (non-equivariant) spectrum functor is the indentity. On the other hand, as a diagram living in the 1-category $\mathrm{CAlg}(\Sp^{\Gr,\heartsuit})$ the underlying diagram is none other than the $\MU$-cobar complex. The algebra ExtMU was defined as the totalization of the $\MU$-cobar complex, so this proves the lemma.
%\end{proof}
%
%\subsection{Homotopy of $C\ta$-modules}
%
%\begin{cnv}
%  Given $M \in \Ab^{\Gr}$, we let $\pi_{p,w} M = \pi_p M_w$.
%  Similarly, given $M \in \uAb^{\Gr}$, we let $\pi_{p,q,w} M = \pi_{p + q \sigma} ^{C_2} M_w$.
%\end{cnv}
%
%\begin{lem} \label{lem:underline-homotopy}
%  Given any abelian group $A$, we have
%  \[\pi_{p+q\sigma} ^{C_2} \underline{A} \cong \pi_{p} \left( \bigoplus_{a} A \otimes \Sigma^a \pi_{a+q\sigma} \uZ \right).\]
%\end{lem}
%
%\begin{proof}
%  Consider the composite
%  \[\Ab \to \uAb \to \Ab^{\Gr}\]
%  given by
%  \[A \mapsto \underline{A} \mapsto \{\Map(\Ss^{q\sigma}, \underline{A})\}\]
%  This composite is colimit-preserving, hence is equivalent to
%  \[A \mapsto \{A \otimes \Map(\Ss^{q\sigma} , \underline{\Z})\}.\]
%  Now, we have
%  \[\Map(\Ss^{q\sigma} , \underline{\Z}) \simeq \bigoplus_a \Sigma^{a} \pi_{a+q\sigma} \uZ \]
%  because it's a $\Z$-module, so the result follows from applying $\pi_p$.
%\end{proof}
%
%\begin{lem}
%  Let $M \in \Ab^{\Gr}$. Then the homotopy groups of $\twist^{\rho} (\underline{M}) \in \uAb^{\Gr}$ are given by
%  \[\pi_{p,q,w} \twist^{\rho} (\underline{M}) \cong \]
%\end{lem}
%
%\begin{proof}
%  It follows from the fact that $\Z$ is a PID that
%  \[M_w \simeq \bigoplus_{p \in \Z} \Sigma^p \pi_{p,w} M.\]
%  It follows that
%  \[\underline{M}_w \simeq \bigoplus_{p \in \Z} \Sigma^p \underline{\pi_{p,w} M},\]
%  and hence that
%  \[(\twist^{\rho} (\underline{M}))_w \simeq \bigoplus_{p \in \Z} \Sigma^{p + w \rho} \underline{\pi_{p,w} M}.\]
%
%  The result now follows from
%\end{proof}
%
%
%
%
%
%Given $M \in \IndCoh(\Mfg)_{i2}$, we set
%\[\Ext_{\MU_* \MU} ^{s,2w} (\MU_*, M) \coloneqq [\Sigma^{s} \omega_{\mathbb{G}/\Mfg} ^{\otimes w}, M].\]
%
%\begin{prop}
%  Let $E \in \IndCoh(\Mfg)_{i2}$, and let
%  \[\underline{E} \in \uAb_{i2} \otimes _{\Ab_{i2}} \IndCoh(\Mfg)_{i2}\]
%  denote the image of $E$ under the twisted underline functor
%  \[\_ : \IndCoh(\Mfg)_{i2} \to \uAb_{i2} \otimes _{\Ab_{i2}} \IndCoh(\Mfg)_{i2},\]
%  and further let
%  \[\underline{E}^{\rho} \in \uAb_{i2} \otimes _{\Ab_{i2}} \IndCoh(\Mfg)_{i2}\]
%  denote the image of $\underline{E}$ under the $\rho$-twist functor
%  \[\twist^{\rho} : \uAb_{i2} \otimes_{\Ab_{i2}} \IndCoh(\Mfg)_{i2} \to \uAb_{i2} \otimes _{\Ab_{i2}} \IndCoh(\Mfg)_{i2},\]
%  Then the homotopy groups of $\underline{E}$, viewed as a $C\ta$-module under the equivalence of categories of \Cref{thm:Cta}, are given by
%  \[ \pi_{p,q,w}^\R \underline{E}^{\rho} \cong \bigoplus_{w+a-s = p} \Ext_{\MU_*\MU}^{s,2w}(\MU_*, E) \otimes_{\Z_2} \pi_{a + (q-w) \sigma}^{C_2} \uZt, \]
%  where the tensor product is implicitly derived.
%\end{prop}
%
%\begin{proof}
%  By [lems], we may view the relevant underline functor as a functor
%  \[\_ : \Mod(\Ab^{\Gr}; \ExtMUt) \to \Mod(\uAb^{\Gr}; \underline{\ExtMUt}),\]
%  and view $\twist^{\rho}$ as a functor
%  \[\twist^{\rho} ; \Mod(\uAb^{\Gr}; \underline{\ExtMUt}) \to \Mod(\uAb^{\Gr}; \underline{\ExtMUt}),\]
%  To compute the homotopy groups, it suffices to compute the image of $\underline{E}^{\rho}$ under the forgetful functor
%  \[\uAb \otimes_{\Ab} \IndCoh(\Mfg)_{i2} \simeq \Mod(\uAb^{\Gr}; \ExtMUt) \to \uAb^{\Gr}.\]
%  This only depends on the image of $E$ under the forgetful functor 
%  \[\IndCoh(\Mfg)_{i2} \simeq \Mod(\Ab^{\Gr}; \ExtMUt) \to \Ab^{\Gr},\]
%  which is $\{\bigoplus_{s} \Sigma^{-s} \Ext_{\MU_* \MU} ^{s,2w} (\MU_*, M)\}_{w \in \Z}.$
%\end{proof}
%


